% ed2_equals_2.tex -- complete proof that ed^[2]_k(tau_gen) = 2 at (p,n) = (2,3).
% Companion to the paper [P] and the addendum [A]; every external input is
% listed in Section 1 and again in Section 8.
\documentclass[11pt]{amsart}
\usepackage[margin=1.1in]{geometry}
\usepackage{amsmath,amssymb,amsthm}
\usepackage{array,booktabs,enumitem}
\usepackage[hidelinks]{hyperref}

\newtheorem{theorem}{Theorem}[section]
\newtheorem{proposition}[theorem]{Proposition}
\newtheorem{lemma}[theorem]{Lemma}
\newtheorem{corollary}[theorem]{Corollary}
\newtheorem{fact}[theorem]{Fact}
\theoremstyle{definition}
\newtheorem{definition}[theorem]{Definition}
\newtheorem{construction}[theorem]{Construction}
\theoremstyle{remark}
\newtheorem{remark}[theorem]{Remark}

\newcommand{\ed}{\operatorname{ed}}
\newcommand{\edd}[1]{\ed^{[#1]}}
\newcommand{\tr}{\operatorname{trdeg}}
\newcommand{\im}{\operatorname{im}}
\newcommand{\gen}{\mathrm{gen}}
\newcommand{\Z}{\mathbb Z}
\newcommand{\A}{\mathbb A}
\newcommand{\PP}{\mathbb P}
\newcommand{\F}{\mathbb F}
\newcommand{\res}{\operatorname{res}}
\newcommand{\cores}{\operatorname{cores}}
\newcommand{\Pic}{\operatorname{Pic}}
\newcommand{\Aut}{\operatorname{Aut}}
\newcommand{\End}{\operatorname{End}}
\newcommand{\Hom}{\operatorname{Hom}}
\newcommand{\Spec}{\operatorname{Spec}}
\newcommand{\Gal}{\operatorname{Gal}}
\newcommand{\Br}{\operatorname{Br}}
\newcommand{\Supp}{\operatorname{Supp}}
\newcommand{\divv}{\operatorname{div}}
\newcommand{\Tr}{\operatorname{Tr}}
\newcommand{\Nm}{\operatorname{N}}
\newcommand{\ut}{\tilde u}
\newcommand{\Ltwo}{L_2}
\newcommand{\Ltwop}{L_2'}
\newcommand{\Kp}{K'}
\newcommand{\Lp}{L'}
\newcommand{\Gtwo}{G_2}
\newcommand{\Etwo}{E_2}
\newcommand{\Ctwo}{C_2}

\setcounter{section}{-1}

\begin{document}

\title[The second slot equals two]{The generic $\Z/8$-torsor in characteristic $2$
stays two-dimensional over every quadratic extension:\\ a complete proof of
$\ed^{[2]}_k(\tau_{\gen})=2$}
\author{}
\date{28 September 2026}
\maketitle

\section{Statement and outline}\label{sec:statement}

Throughout, $k$ is an algebraically closed field of characteristic $2$,
$G=\Z/8=\langle g\rangle$, $L=k(z_0,z_1,z_2)$ with $G$ acting by the Witt translation
$z\mapsto z+1$, $K=L^G=k(s_0,s_1,s_2)$ where $s=(s_0,s_1,s_2)=\wp(z)$, and
$\tau_{\gen}\colon\Spec L\to\Spec K$ is the generic $\Z/8$-torsor, of class
$s\in W_3(K)/\wp W_3(K)=H^1(K,\Z/8)$. For a $G$-torsor $\tau$ over a field $F\supseteq k$,
$\ed_k(\tau)$ is the least transcendence degree over $k$ of a subfield $k\subseteq F_0\subseteq F$
such that $\tau\cong\tau_0\otimes_{F_0}F$ for a $G$-torsor $\tau_0$ over $F_0$
(\cite[\S1]{BF}, \cite[Def.~3.1]{P}), and
\[
\edd d_k(\tau):=\min\{\ed_k(\tau\otimes_FF'):[F':F]\le d\},\qquad
d_j(\tau):=\min\{d\ge1:\edd d_k(\tau)\le j\}
\]
(\cite[Definition 3.1]{P}). The paper \cite{P} proved $\edd2_k(\tau_{\gen})\in\{1,2\}$
and $d_1(\tau_{\gen})\in\{2,3,4\}$ \cite[Theorem 7.1]{P}. The purpose of this note is to
settle the second slot.

\begin{theorem}\label{thm:main}
Let $k$ be algebraically closed of characteristic $2$, $K=k(s_0,s_1,s_2)$ and
$\tau_{\gen}$ the generic $\Z/8$-torsor of Witt class $s=(s_0,s_1,s_2)$. Then
\[
\ed_k(\tau_{\gen}\otimes_KK')\ \ge\ 2\qquad\text{for every field extension $K'/K$ of degree $2$,}
\]
with equality for $K'=K(v_2)$, $v_2^2+v_2=s_2$. Consequently
\[
\edd2_k(\tau_{\gen})=2\qquad\text{and}\qquad d_1(\tau_{\gen})\in\{3,4\}.
\]
\end{theorem}

\begin{remark}\label{rem:statement}
(i) Since $\ed_k(\tau_{\gen}\otimes K')\le\ed_k(\tau_{\gen})\le3$, the value
$\ed_k(\tau_{\gen}\otimes K')$ lies in $\{2,3\}$ for every quadratic $K'$, and it equals
$2$ for all the extensions listed in \cite[Theorem 7.1]{P} and \cite[Theorem 11.9]{A};
it equals $2$ for every quadratic $K'$ as soon as $\ed_k(\Z/8)=2$. We do not claim the
upper bound $2$ for an arbitrary quadratic $K'$: what is proved here, and what
$\edd2_k(\tau_{\gen})=2$ means, is the lower bound.
(ii) The theorem is stated for $k$ algebraically closed because every input from
\cite{P} is.
\end{remark}

\subsection*{Outline of the proof}
Suppose $\ed_k(\tau_{\gen}\otimes K')\le1$ for some quadratic $K'/K$. By purely
inseparable invariance and the exclusion of $K(z_0)$, $K'=K(v_f)$ is separable with
$f\notin\{0,s_0\}+\wp K$, so $L'=L\otimes_KK'$ is a field. The geometric criterion of
\cite[Theorem 4.2, Corollary 4.4]{P} turns the hypothesis into a smooth projective curve
$C$ over $k$ with faithful $G$-action and a $G$-equivariant field embedding
$\iota\colon k(C)\hookrightarrow L'$; the fixed field $k(B)$, $B=C/G$, embeds into $K'$
and is the ``descent field'' $F_0$. We pass to the quotient $C_2=C/\langle g^4\rangle$
with its faithful $\Z/4$-action and to its twist $E_2$ by the $\Z/4$-class $(s_0,s_1)_K$;
the embedding $\iota$ produces a point $Q\in E_2(K')$ whose image in $B$ is the descent
point, hence is non-constant. The $K$-rational divisor $D_Q$ of degree $2$ on $E_2$
attached to $Q$ (``$Q+\sigma Q$'') is the heart of the argument: after the base change
$L_2=k(z_0,z_1,s_2)$ that trivialises the twist, its class is that of a
\emph{constant} divisor $D_0$ on $C_2$ with $\Z/4$-invariant class, because rational
maps from $\A^3$ to a Jacobian are constant (Step~1). Either $\ell(D_0)=1$, and then
$D_Q$ is supported on constant points, forcing the descent point to be constant
(Step~2); or $\ell(D_0)=2$, and then $D_0$ moves in a pencil $|D_0|$ that is the pullback
of a degree-$2$ map $u\colon C_2\to\PP^1$ on which $\Z/4$ acts through $u\mapsto u+1$
(the hyperelliptic case, Step~3, or the supersingular genus-$1$ case, Step~4(i)); the
twisted coordinate $\ut=u+z_0$ is $K$-rational, $D_Q$ is a $\ut$-fibre over $\PP^1(K)$,
and the descent point $t(Q)=\ut(Q)^2+\ut(Q)+s_0$ lies in $K$. That puts the descent
field inside $K$, which is impossible by the truncation argument of
\cite[Theorem 11.6(b)]{A}, reproved here. The only remaining case is an ordinary elliptic
curve $C_2$ on which $\Z/4$ acts by translation (Step~4(ii)); there the class
$(s_0,s_1)_K\otimes K'$ is the isogeny torsor at a point $b\in B(K')$, and the
corestriction identity $\cores\circ\delta_{K'}=\delta_K\circ\Nm$ together with
$B(K)=B(k)$ gives $2\,(s_0,s_1)=0$ in $H^1(K,\Z/4)$, which is false. Section~7 checks
that the cases are exhaustive and derives the corollaries; Section~8 lists what is cited
without proof.

\section{Conventions and imported facts}\label{sec:facts}

We use the notation of \cite[Sections 2--7]{P} and \cite[Section 11]{A}; [P] and [A]
are cited with their own numbering. All curves are over $k$ unless said otherwise; a
\emph{curve over a field $F$} is a smooth projective geometrically integral $F$-scheme
of dimension $1$. For a divisor $D$ on a curve $X$ over $F$ we write
$L(D)=H^0(X,\mathcal O_X(D))$ and $\ell(D)=\dim_FL(D)$. Fix
\[
G_2:=G/\langle g^4\rangle\cong\Z/4,\qquad K_2:=k(s_0,s_1),\qquad
L_2:=k(z_0,z_1,s_2)=L^{\langle g^4\rangle}.
\]
The generator $g$ of $G$ and its image in $G_2$ are both written $g$.

\subsection{Witt vectors}\label{ss:witt}
$W_n(A)=A^n$ with the ring structure defined by the ghost components
$w_m(x)=\sum_{i\le m}p^ix_i^{p^{m-i}}$; $F(x)=(x_i^p)$ on $\F_p$-algebras,
$V(x)=(0,x_0,\dots,x_{n-2})$, $\wp=F-\mathrm{id}$, $[a]=(a,0,\dots,0)$
\cite[\S2.2]{P}. We record the formulas we use, at $p=2$.

\begin{fact}[Witt arithmetic at $p=2$]\label{fact:witt}
\begin{enumerate}[label=\textup{(\alph*)}]
\item $(a+b)_0=a_0+b_0$, $(a+b)_1=a_1+b_1+a_0b_0$, and
$(a+b)_2=a_2+b_2+a_1b_1+a_0b_0(a_1+b_1)+a_0^3b_0+a_0b_0^3$ \cite[(3)]{P}.
\item $\wp(a_0,a_1)=(a_0^2+a_0,\;a_1^2+a_1+a_0^3+a_0^2)$ and
$\wp(a)_2=a_2^2+a_2+C(a_0,a_1)$ with $C$ as in \cite[(4)]{P}.
\item $x+V^r(y)=(x_0,\dots,x_{r-1},y_0,\dots)$ if $x_r=\dots=x_{n-1}=0$
\cite[(2)]{P}; $V$ is additive and $\wp V=V\wp$.
\item $2\cdot(a_0,a_1)=(0,a_0^2)$ in $W_2(A)$ and $2\cdot(a_0,a_1,a_2)=(0,a_0^2,a_1^2)$
in $W_3(A)$, for every $\F_2$-algebra $A$.
\item $(0,c)\in\wp W_2(F)$ iff $c\in\wp F$; $(0,c_0,c_1)\in\wp W_3(F)$ iff
$(c_0,c_1)\in\wp W_2(F)$ \cite[(5)]{P}.
\end{enumerate}
\end{fact}

\begin{proof}
(a),(b),(c) are \cite[\S2.2]{P}; for the reader's convenience: $w_1(a+b)=w_1(a)+w_1(b)$
reads $(a_0+b_0)^2+2(a+b)_1=a_0^2+2a_1+b_0^2+2b_1$, so $(a+b)_1=a_1+b_1-a_0b_0$
as a polynomial identity over $\Z$ (the ghost map is injective on $W_2(\Z[a,b])$), and
reduction modulo $2$ gives (a); (b) is $F-\mathrm{id}$ with $F$ the componentwise
square. (d) By (a) with $a=b$: $(2a)_0=2a_0=0$ and $(2a)_1=2a_1+a_0^2=a_0^2$; in
$W_3$, $(2a)_2=2a_2+a_1^2+a_0^2\cdot2a_1+2a_0^4=a_1^2$. (e) If $(0,c)=\wp(x_0,x_1)$
then $x_0^2+x_0=0$, so $x_0\in\{0,1\}$, $x_0^3+x_0^2=0$, and $c=x_1^2+x_1\in\wp F$;
conversely $(0,\wp x_1)=\wp(0,x_1)$ by (b). For $W_3$: if $(0,c_0,c_1)=\wp(x)$ then
again $x_0\in\{0,1\}$, $\wp[x_0]=0$, and $x=[x_0]+V(x_1,x_2)$ gives
$V(c_0,c_1)=\wp V(x_1,x_2)=V\wp(x_1,x_2)$, so $(c_0,c_1)=\wp(x_1,x_2)$ as $V$ is
injective on components; the converse is $\wp V=V\wp$.
\end{proof}

\begin{fact}[Artin--Schreier--Witt theory]\label{fact:asw}
For every field $F$ of characteristic $2$ there is a group isomorphism
$H^1(F,\Z/2^n)\cong W_n(F)/\wp W_n(F)$, functorial in $F$, sending the class of $a$ to
the torsor $\Spec F[x]/(\wp x=a)$ with $1\in\Z/2^n$ acting by $x\mapsto x+1$.
Truncation $W_n\to W_m$ corresponds to $\Z/2^n\to\Z/2^m$ and $V^{n-m}\colon W_m\to W_n$
to the inclusion $\Z/2^m\hookrightarrow\Z/2^n$; both commute with $\wp$ and with base
change. \cite[\S2.2]{P}, \cite[Ch.~V]{Se}.
\end{fact}

\begin{lemma}\label{lem:Vinj}
$V\colon W_2(F)/\wp\to W_3(F)/\wp$ is injective, i.e.\ $H^1(F,\Z/4)\to H^1(F,\Z/8)$ is
injective. Moreover $s_0\notin\wp K$, hence $(0,s_0)\ne0$ in $W_2(K)/\wp$ and
$(0,s_0,s_1)\ne0$ in $W_3(K)/\wp$; in particular $(s_0,s_1)\ne0$ in $W_2(K)/\wp$.
\end{lemma}

\begin{proof}
Injectivity is Fact~\ref{fact:witt}(e). If $s_0=a^2+a$ with $a\in K=k(s_1,s_2)(s_0)$,
then $a$ is integral over $k(s_1,s_2)[s_0]$, which is integrally closed, so $a$ is a
polynomial in $s_0$ of some degree $d$, and $a^2+a$ has degree $2d\ne1$: contradiction.
The rest follows from Fact~\ref{fact:witt}(e), and $(s_0,s_1)\ne0$ because its
truncation $s_0$ is nonzero in $K/\wp K$.
\end{proof}

\subsection{The generic torsor and its truncations}
By \cite[\S2.3, Proposition 3.9]{P}, $L/K$ is Galois with group $G$, its unique
quadratic subextension is $K_1=L^{\langle g^2\rangle}=K(z_0)$ (with $z_0^2+z_0=s_0$), and
$L^{\langle g^4\rangle}=L_2=k(z_0,z_1,s_2)$; the group $G_2=G/\langle g^4\rangle$ acts on $L_2$
by the Witt translation $(z_0,z_1)\mapsto(z_0,z_1)+(1,0)$ of $W_2$, and
$\wp(z_0,z_1)=(s_0,s_1)$ (truncation of $\wp z=s$). Hence
\[
T_2:=\Spec L_2\longrightarrow\Spec K
\]
is the $\Z/4$-torsor of class $(s_0,s_1)\in W_2(K)/\wp$, written $(s_0,s_1)_K$; it is
$\tau_2\otimes_{K_2}K$ where $\tau_2$ is the generic $\Z/4$-torsor over $K_2=k(s_0,s_1)$.

\begin{fact}\label{fact:ed}
\begin{enumerate}[label=\textup{(\alph*)}]
\item \cite[Lemma 2.1]{P} $\ed_k(\tau_{\gen})=\ed_k(\Z/8)\le3$; \cite[Prop.~6.1]{P}
$\ed_k(\tau_2)=\ed_k(\Z/4)=2$.
\item \cite[Lemma 4.1(i)]{P} If $\tau$ is a $G$-torsor over an infinite field
$F\supseteq k$ and $t$ is transcendental over $F$, then $\ed_k(\tau\otimes F(t))=\ed_k(\tau)$.
Consequently $\ed_k((s_0,s_1)_K)=\ed_k(\tau_2\otimes_{K_2}K_2(s_2))=2$.
\item \cite[Prop.~3.2(c)]{P} $\edd d_k(\tau_{\gen})\ge1$ for $d<8$.
\item \cite[Prop.~3.6]{P} If $F'/F$ is finite purely inseparable, then
$\ed_k(\tau\otimes F')=\ed_k(\tau)$; the profile only sees separable extensions.
\item \cite[Prop.~3.9]{P} $\ed_k(\tau_{\gen}\otimes K(z_0))\ge\min(n-1,2)=2$.
\item \cite[Theorem 7.1]{P} $\edd2_k(\tau_{\gen})\le2$, witnessed by $K(v_2)$,
$v_2^2+v_2=s_2$ (class $(s_0,s_1,0)$, defined over $k(s_0,s_1)$); $\edd4_k(\tau_{\gen})=1$;
$\ed_k(\tau_{\gen})\in\{2,3\}$.
\item \cite[\S3.1]{P}, \cite[Lemma 2.1(b)]{BF} $\ed_k(\tau\otimes F')\le\ed_k(\tau)$ for
every extension $F'/F$.
\end{enumerate}
\end{fact}

\subsection{Twists and the geometric criterion}
For a quasi-projective $G$-variety $X$ over $k$ and a $G$-torsor $\tau\colon T\to\Spec F$,
the twist is ${}^\tau X=(T\times_kX)/G$, an $F$-variety; it commutes with base change in
$F$ \cite[\S2.4]{P}, \cite[\S3]{DR}.

\begin{fact}[{\cite[Fact 2.2(a)]{P}, \cite[Lemma 7.3]{DR}}]\label{fact:twistpts}
For every field $F'\supseteq F$, ${}^\tau X(F')=\Hom_G(T\times_FF',X)$, the set of
$G$-equivariant $k$-morphisms.
\end{fact}

\begin{fact}[{\cite[Theorem 4.2, (1)$\Rightarrow$(3) and last paragraph]{P}}]\label{fact:geom}
Let $\tau\colon T\to\Spec F$ be a $G$-torsor and suppose $\ed_k(\tau\otimes F')\le m$ for
a field $F'/F$. Then there is a quasi-projective generically free primitive $G$-variety
$Z$ over $k$ with $\dim Z\le m$ and a $G$-equivariant morphism $T\times_FF'\to Z^{\rm free}$,
where $Z^{\rm free}=Z\smallsetminus Z^{g^4}$ is the locus of points with trivial
stabiliser. Conversely such a map gives $\ed_k(\tau\otimes F')\le\dim Z$ \cite[Lemma 2.3]{P}.
If $T\times_FF'$ is the spectrum of a field, $Z$ may be taken irreducible with faithful
$G$-action.
\end{fact}

\begin{fact}[{\cite[Lemma 2.4]{P}}]\label{fact:image}
Let $Z$ be an irreducible $G$-variety with faithful action, $T$ connected, and
$f\colon T\to Z$ equivariant. The closure $Y$ of the image point of $f$ is irreducible and
$G$-stable, and the image point is free iff $G$ acts faithfully on $Y$. If $\dim Z=1$
and the image point is free, then $f$ is dominant.
\end{fact}

\begin{fact}[{\cite[Lemma 2.5(ii)]{P}}]\label{fact:pgl}
$\mathrm{PGL}_2(k)$ has no element of order $4$. Every element of order $2$ of
$\mathrm{PGL}_2(k)$ is conjugate to $u\mapsto u+1$.
\end{fact}

\begin{proof}
A $2$-element of $\mathrm{PGL}_2(k)$ lifts to $A\in\mathrm{GL}_2(k)$ with $A^{2^m}=\lambda$
scalar; replacing $A$ by $A/\mu$ with $\mu^{2^m}=\lambda$ ($k$ is algebraically closed) we
may assume $A^{2^m}=1$.
Then $(A-1)^{2^m}=A^{2^m}-1=0$, so $N:=A-1$ is nilpotent, $N^2=0$, and
$A^2=1+2N=1$: no element of order $4$. If the image of $A$ has order $2$ then $N\ne0$,
so $N$ is conjugate to $\left(\begin{smallmatrix}0&1\\0&0\end{smallmatrix}\right)$ and $A$
to $\left(\begin{smallmatrix}1&1\\0&1\end{smallmatrix}\right)$, which acts on $\PP^1$ by
$u\mapsto u+1$.
\end{proof}

\subsection{Abelian varieties, rational maps, Galois descent}

\begin{fact}[rigidity]\label{fact:rigid}
Let $A$ be an abelian variety over $k$, and let $P$ be a $k$-variety that is a torsor
under $A$ (a principal homogeneous space).
\begin{enumerate}[label=\textup{(\alph*)}]
\item Every $k$-morphism $\PP^1_k\to A$ is constant \cite[Cor.~3.8]{Mi}; equivalently
\cite[Lemma 2.5(iv)]{P}.
\item (Weil) Every rational map from a smooth $k$-variety to $A$ extends to a morphism
\cite[Thm.~3.1]{Mi}, \cite[8.4, Cor.~6]{BLR}.
\item Every $k$-morphism $\Spec F\to P$, where $F=k(t_1,\dots,t_m)$ is purely
transcendental over $k$, factors through a $k$-point of $P$.
\item \cite[Lemma 4.1(ii)]{P} If $F\supseteq k$ is a field and $F'=F(t_1,\dots,t_m)$ is
purely transcendental over $F$, then $A(F')=A(F)$; the same holds for an abelian variety
over $F$.
\end{enumerate}
\end{fact}

\begin{proof}
(a),(b) are cited. (c) Since $k$ is algebraically closed, $P(k)\ne\emptyset$ and $P\cong A$
as $k$-varieties (translate a $k$-point to the origin); so a $k$-morphism
$\Spec F\to P$ is a rational map $\A^m_k\dashrightarrow A$, which by (b) extends to a
morphism $\phi\colon\A^m\to A$. For $x\in\A^m(k)$ let $\ell_x\subseteq\A^m$ be the line
through $0$ and $x$ (a point if $x=0$); $\phi|_{\ell_x}$ extends to $\PP^1\to A$ and is
constant by (a), so $\phi(x)=\phi(0)$ for all $x$, and $\phi$ is constant ($\A^m$ is
reduced and $k$ algebraically closed). Hence $\Spec F\to P$ factors through the
$k$-point $\phi(0)$. (d) The proof in \cite{P} is written for an elliptic curve but uses
only that $A$ is proper and that $\PP^1_{\bar F}\to A_{\bar F}$ is constant (a); it
applies verbatim to any abelian variety over $F$.
\end{proof}

\begin{fact}[Galois descent]\label{fact:descent}
Let $M/F$ be a finite Galois extension with group $\Gamma$.
\begin{enumerate}[label=\textup{(\alph*)}]
\item For $F$-schemes $S,X$: $\Hom_F(S,X)\to\Hom_M(S_M,X_M)$ is injective with image
the morphisms commuting with the natural $\Gamma$-actions; in particular
$X(M)^\Gamma=X(F)$ \cite[6.1, Thm.~6 and 6.2, Example B]{BLR}.
\item A descent datum on a quasi-projective $M$-scheme $Y$ relative to $M/F$, i.e.\ a
$\Gamma$-action by automorphisms $\phi_\gamma$ of the $F$-scheme $Y$ covering
$\Spec(\gamma)$ on $\Spec M$, is effective: there is an $F$-scheme $X$, unique up to
unique isomorphism, with an $M$-isomorphism $\psi\colon X_M\to Y$ such that
$\psi\circ(\mathrm{id}_X\times\Spec\gamma)=\phi_\gamma\circ\psi$ for all $\gamma$
\cite[6.2, Example B]{BLR}, \cite[V.20, Prop.~12]{SeAG}. $X$ is a curve over $F$ iff $Y$
is a curve over $M$ (these properties are local for the fpqc topology). $F$-morphisms
$X\to X'$ between descended schemes are the $M$-morphisms $Y\to Y'$ compatible with the
descent data (by (a)). If $X$ is a curve then $F(X)=M(Y)^\Gamma$, where $\Gamma$ acts on
$M(Y)=M(X_M)=F(X)\otimes_FM$ through the $\phi_\gamma^*$.
\item If $X$ is an $F$-scheme, $X_M$ carries the descent datum $\mathrm{id}_X\times\Spec\gamma$,
and (a) is the special case of (b) for this datum. A Galois-invariant divisor on $X_M$
descends uniquely to a divisor on $X$ (for $X$ a curve: closed points of $X$ correspond
to $\Gamma$-orbits of closed points of $X_M$, with $x\otimes_FM=\sum_{y\mapsto x}y$).
\end{enumerate}
\end{fact}

\begin{fact}[Picard schemes]\label{fact:pic}
Let $C$ be a curve over $k$ of genus $g$.
\begin{enumerate}[label=\textup{(\alph*)}]
\item The Picard functor of $C/k$ is represented by a smooth $k$-group scheme
$\Pic_{C/k}$, locally of finite type, whose identity component $J=\Pic^0_{C/k}$ is an
abelian variety of dimension $g$ (the Jacobian) and whose connected components are the
$\Pic^d_{C/k}$, $d\in\Z$, each a torsor under $J$ \cite[8.2, Thm.~3; 8.4, Thm.~3]{BLR},
\cite[Thm.~9.4.8, Prop.~9.5.3, Cor.~9.5.4]{Kl}.
\item Since $C(k)\ne\emptyset$, for every field $F\supseteq k$ the natural map
$\Pic(C_F)\to\Pic_{C/k}(F)$ is bijective and respects degrees
\cite[8.1, Prop.~4]{BLR}. Under it, for $\alpha\in\Aut(C)$ the pullback $\alpha_F^*$ on
$\Pic(C_F)$ corresponds to $(\alpha^*)_F$, where $\alpha^*$ is the automorphism of the
$k$-group scheme $\Pic_{C/k}$ induced by functoriality.
\end{enumerate}
\end{fact}

\begin{fact}[cohomology and base change]\label{fact:bc}
Let $X$ be a curve over $F$, $F'/F$ a field extension, $D$ a divisor on $X$.
\begin{enumerate}[label=\textup{(\alph*)}]
\item $L(D_{F'})=L(D)\otimes_FF'$; in particular $\ell(D_{F'})=\ell(D)$ and
$\deg D_{F'}=\deg D$ \cite[III.9.3]{Ha}.
\item $\Pic(X)\to\Pic(X_{F'})$ is injective.
\item Two $F$-morphisms $X\to Y$ that agree after base change to $F'$ are equal.
\item For a nonconstant $h\in F(X)$ and the associated morphism $h\colon X\to\PP^1_F$,
the fibre divisor $h^*(c)$, $c\in\PP^1(F)$, is effective, $F$-rational, of degree
$\deg h=[F(X):F(h)]$, and $h^*(c)_{F'}=(h_{F'})^*(c)$.
\end{enumerate}
\end{fact}

\begin{proof}
(a) is flat base change. (b) If $\deg D=0$ and $D_{F'}$ is principal then
$\ell(D)=\ell(D_{F'})=1$ by (a); a nonzero $h\in L(D)$ has $\divv h+D\ge0$ of degree $0$,
so $D=\divv(h^{-1})$. (c) $\Spec F'\to\Spec F$ is faithfully flat, hence an epimorphism
of schemes. (d) is standard \cite[II.6.9, II.6.10]{Ha}; pullback of Cartier divisors
commutes with base change.
\end{proof}

\subsection{Curves: linear systems, hyperelliptic curves, elliptic curves}

\begin{fact}[Riemann--Roch and consequences]\label{fact:rr}
Let $C$ be a curve over $k$ of genus $g$ and $D$ a divisor on $C$.
\begin{enumerate}[label=\textup{(\alph*)}]
\item $\ell(D)-\ell(K_C-D)=\deg D+1-g$ \cite[IV.1.3]{Ha}. If $g=1$ and $\deg D>0$
then $\ell(D)=\deg D$.
\item If $g\ge1$ and $\deg D\ge1$ then $\ell(D)\le\deg D$.
\item If $g\ge1$, $\deg D=2$ and $\ell(D)=2$, then $|D|$ is base point free and defines a
morphism $u\colon C\to\PP^1$ of degree $2$ with $D\sim u^*(a)$ for every $a\in\PP^1(k)$;
$u$ is separable, and $k(C)/k(u)$ is Galois with group $\{1,\iota\}$ generated by an
involution $\iota$ of $C$ with $C/\iota\cong\PP^1$. Conversely the fibres of any
degree-$2$ morphism $C\to\PP^1$ form such a pencil.
\item (uniqueness of the $g^1_2$) If $g\ge2$, any two divisors $D,D'$ of degree $2$ with
$\ell(D)=\ell(D')=2$ are linearly equivalent.
\end{enumerate}
\end{fact}

\begin{proof}
(a) is cited; for $g=1$, $K_C=0$ and $\ell(-D)=0$. (b) For a point $P$,
$\ell(D)\le\ell(D-P)+1$ (evaluation at $P$ in a local coordinate). By induction on
$\deg D$ this gives $\ell(D)\le\ell(D-(\deg D-1)P)+\deg D-1$ with $\deg(D-(\deg D-1)P)=1$;
and a divisor $D_1$ of degree $1$ with $\ell(D_1)\ge2$ gives a nonconstant $h$ with
$\divv h+D_1\ge0$, hence a degree-$1$ map $C\to\PP^1$, i.e.\ $C\cong\PP^1$ and $g=0$. So
$\ell(D_1)\le1$ and $\ell(D)\le\deg D$. (c) If $P$ were a base point of $|D|$ then
$\ell(D-P)=\ell(D)=2$ with $\deg(D-P)=1$, contradicting (b). So $|D|$ is base point
free; with $L(D)=\langle h_0,h_1\rangle$ the map $u=h_1/h_0\colon C\to\PP^1$ has
$u^*(a)=\divv(h_1-ah_0)+D\sim D$ (and $u^*(\infty)=\divv(h_0)+D$), and $\deg u=\deg D=2$.
If $k(C)/k(u)$ were inseparable of degree $2$, it would be purely inseparable, so
$k(C)=k(u^{1/2})$ and $C\cong\PP^1$, contradicting $g\ge1$; a separable extension of
degree $2$ is Galois, and the nontrivial element of its Galois group is induced by an
involution $\iota$ of $C$ with $k(C/\iota)=k(C)^\iota=k(u)$. The converse is clear.
(d) Let $u,u'\colon C\to\PP^1$ be the degree-$2$ maps of (c) for $D,D'$, and let
$\Gamma\subseteq\PP^1\times\PP^1$ be the image of $(u,u')$, an irreducible curve of
bidegree $(a,b)$, and $e:=\deg(C\to\Gamma)$. The projections give $eb=\deg u=2$ and
$ea=\deg u'=2$. If $e=2$ then $a=b=1$ and $\Gamma$ is the graph of an automorphism
$\alpha$ of $\PP^1$, so $u'=\alpha\circ u$ and $D'\sim u'^*(a_0)=u^*(\alpha^{-1}a_0)\sim D$.
If $e=1$ then $\Gamma$ has bidegree $(2,2)$, arithmetic genus $(2-1)(2-1)=1$, and $C$ is
its normalisation, so $g\le1$: excluded.
\end{proof}

\begin{fact}[elliptic curves in characteristic $2$]\label{fact:ell}
Let $E$ be an elliptic curve over $k$ with origin $O$.
\begin{enumerate}[label=\textup{(\alph*)}]
\item $\Aut(E)=E(k)\rtimes\Aut(E,O)$: every automorphism of the curve $E$ is uniquely
$t_P\circ\alpha$ with $P\in E(k)$, $t_P$ the translation, $\alpha\in\Aut(E,O)$, and
every $\alpha\in\Aut(E,O)$ is a group automorphism, $\alpha\circ t_P=t_{\alpha P}\circ\alpha$
\cite[III.4.8]{Sil}.
\item $E$ is supersingular iff $j(E)=0$ iff $E(k)[2]=0$; if $E$ is ordinary then
$E(k)[2^r]\cong\Z/2^r$ for all $r$ \cite[V.3.1, Exercise 5.7]{Sil}. If $j(E)\ne0$ then
$\Aut(E,O)=\{\pm1\}$; if $j(E)=0$ then $\Aut(E,O)$ has order $24$ with Sylow
$2$-subgroups isomorphic to the quaternion group $Q_8$ \cite[App.~A, Prop.~1.2]{Sil},
\cite[Lemma 2.5(iii)]{P}.
\item $\End(E)$ is a domain, torsion-free as a $\Z$-module; for $\phi\in\End(E)$,
$\phi+\hat\phi=1+\deg\phi-\deg(1-\phi)\in\Z$ and $\phi^2-(\phi+\hat\phi)\phi+\deg\phi=0$
\cite[III.4.2, III.6.2, III.6.3]{Sil}. A nonzero isogeny is surjective on $k$-points,
and $\#\ker(\phi)(k)\le\deg\phi$ \cite[II.2.3, III.4.10]{Sil}.
\item For a finite subgroup $\Phi\subseteq E(k)$ there is an elliptic curve $E'$ and a
separable isogeny $\phi\colon E\to E'$ with kernel $\Phi$, unique up to isomorphism; if
$\Phi$ is cyclic of order $4$ generated by $P$, then $\phi$ is finite \'etale of degree
$4$ and $E'\cong E/\langle t_P\rangle$ \cite[III.4.12]{Sil}.
\end{enumerate}
\end{fact}

\subsection{Galois cohomology}

\begin{fact}[restriction, corestriction, connecting maps]\label{fact:gc}
Let $F'/F$ be a finite separable extension of fields, $\Gamma_F\supseteq\Gamma_{F'}$
the absolute Galois groups (for a separable closure $F^s\supseteq F'$), and $M$ a
discrete $\Gamma_F$-module.
\begin{enumerate}[label=\textup{(\alph*)}]
\item $\res\colon H^q(F,M)\to H^q(F',M)$ and $\cores\colon H^q(F',M)\to H^q(F,M)$ are
morphisms of cohomological $\delta$-functors; in degree $0$, $\cores$ is the norm
$M^{\Gamma_{F'}}\to M^{\Gamma_F}$, $m\mapsto\sum_{\gamma\in\Gamma_F/\Gamma_{F'}}\gamma m$;
and $\cores\circ\res=[F':F]$ \cite[Prop.~1.5.2, Cor.~1.5.7]{NSW}, \cite[I.2.4]{Se}.
\item For a finite abelian group $A$ with trivial action, $H^1(F,A)$ is the group of
isomorphism classes of $A$-torsors over $F$ (finite \'etale $F$-schemes with a simply
transitive $A$-action) \cite[I.5.2]{Se}; $\res$ is base change; a torsor that has an
$F$-point is trivial; an $A$-equivariant $F$-morphism between $A$-torsors over $F$ is an
isomorphism \cite[proof of Lemma 2.3]{P}. For $A=\Z/2^n$ this is the group of
Fact~\ref{fact:asw}.
\item Let $0\to A\to E\xrightarrow{\phi}B\to0$ be an exact sequence of commutative
group varieties over $k$ with $\phi$ finite \'etale and $A$ finite constant. Then
$0\to A\to E(F^s)\to B(F^s)\to0$ is an exact sequence of $\Gamma_F$-modules, and the
connecting homomorphism $\delta_F\colon B(F)\to H^1(F,A)$ sends $b$ to the class of the
$A$-torsor $\phi^{-1}(b)=E\times_{B,b}\Spec F$ (with $A$ acting by translation)
\cite[I.5.4--5.6]{Se}. By (a), $\cores_{F'/F}\circ\delta_{F'}=\delta_F\circ\Nm_{F'/F}$
on $B(F')$, where $\Nm_{F'/F}(b)=\sum_{\gamma\in\Gal(F'/F)}\gamma b\in B(F)$ when $F'/F$
is Galois.
\end{enumerate}
\end{fact}

\begin{proof}
Only (c) needs comment. $E(F^s)\to B(F^s)$ is surjective because $\phi^{-1}(b)$ is a
finite \'etale $F^s$-scheme, hence a finite disjoint union of copies of $\Spec F^s$. The
description of $\delta_F(b)$: choose $e\in E(F^s)$ with $\phi(e)=b$; the cocycle
$\gamma\mapsto\gamma(e)-e\in A$ is by definition $\delta_F(b)$, and it is the cocycle of
the torsor $\phi^{-1}(b)$ with respect to its $F^s$-point $e$. The last identity is (a)
applied to the sequence of $\Gamma_F$-modules, restricted to $\Gamma_{F'}$: $\cores$
commutes with $\delta$, and on $H^0$ it is the norm; the sum over $\Gamma_F/\Gamma_{F'}$
is the sum over $\Gal(F'/F)$ when $F'/F$ is Galois.
\end{proof}

\begin{fact}[{\cite[Theorem 11.6]{A}}]\label{fact:caseA}
The statement is reproved below as Theorem~\ref{thm:caseA}; the only inputs are
Facts~\ref{fact:asw}, \ref{fact:ed}(a),(b) and inflation--restriction.
\end{fact}

\section{Setup: a hypothetical witness and the point \texorpdfstring{$Q\in E_2(K')$}{Q in E2(K')}}\label{sec:setup}

\subsection{Reduction to a separable witness with \texorpdfstring{$L\otimes_KK'$}{L(x)K'} a field}

\begin{definition}\label{def:witness}
A \emph{witness} is a separable quadratic extension $K'/K$ with
$\ed_k(\tau_{\gen}\otimes_KK')=1$.
\end{definition}

\begin{lemma}\label{lem:reduce}
If $\ed_k(\tau_{\gen}\otimes_KF')\le1$ for some field extension $F'/K$ of degree $2$,
then $F'=K(v_f)$ with $v_f^2+v_f=f\in K\smallsetminus\wp K$ is a witness and
$f\notin\{0,s_0\}+\wp K$, i.e.\ $F'\ne K(z_0)$.
\end{lemma}

\begin{proof}
If $F'/K$ were purely inseparable, Fact~\ref{fact:ed}(d),(f) would give
$\ed_k(\tau_{\gen}\otimes F')=\ed_k(\tau_{\gen})\ge2$. So $F'/K$ is separable of degree
$2$, hence an Artin--Schreier extension $K(v_f)$, $f\notin\wp K$. By Fact~\ref{fact:ed}(c),
$\ed_k(\tau_{\gen}\otimes F')\ge1$, so it equals $1$. By Fact~\ref{fact:ed}(e), $F'\ne K(z_0)$;
since $z_0^2+z_0=s_0$, $K(z_0)=K(v_{s_0})$, and two Artin--Schreier extensions $K(v_f)$,
$K(v_{f'})$ coincide iff $f\equiv f'$ mod $\wp K$ (Fact~\ref{fact:asw} for $n=1$), so
$f\not\equiv s_0$; and $f\not\equiv0$ as $[F':K]=2$.
\end{proof}

From now on $K'=K(v_f)$ is a witness, $\sigma$ denotes the generator of
$\Gal(K'/K)$, and $f\notin\{0,s_0\}+\wp K$.

\begin{lemma}\label{lem:Lprime}
Put $L':=L\otimes_KK'$ and $L_2':=L_2\otimes_KK'$.
\begin{enumerate}[label=\textup{(\alph*)}]
\item $L'$ is a field, Galois over $K'$ with group $G$ (acting through the first factor),
and $(L')^G=K'$. The involution $\sigma$ acts on $L'$ through the second factor and
commutes with $G$.
\item $L_2'=(L')^{\langle g^4\rangle}$ is a field, Galois over $K'$ with group $G_2$, and
$\Spec L_2'\to\Spec K'$ is the base change $T_2\times_KK'$, i.e.\ the $\Z/4$-torsor of
class $\res_{K'/K}(s_0,s_1)=(s_0,s_1)_K\otimes K'$.
\end{enumerate}
\end{lemma}

\begin{proof}
(a) $L'=L[v]/(v^2+v-f)$, which is a field iff $f\notin\wp L$. By inflation--restriction
for $L/K$ (Galois with group $G$, trivial action on $\Z/2$),
$\ker(H^1(K,\Z/2)\to H^1(L,\Z/2))=H^1(G,\Z/2)=\Hom(\Z/8,\Z/2)\cong\Z/2$, generated by the
class of the unique quadratic subextension $K(z_0)$ of $L/K$, which is $[s_0]$. Hence
$f\in\wp L$ iff $f\in\{0,s_0\}+\wp K$, which is excluded. Thus $L'$ is a field; it is the
base change of the Galois extension $L/K$, so it is Galois over $K'$ with group $G$ and
fixed field $K\otimes_KK'=K'$. The statement about $\sigma$ is clear.
(b) $L_2\otimes_KK'\subseteq L\otimes_KK'$ is a $K'$-subalgebra of a field, finite over
$K'$, hence a field, of degree $[L_2:K]=4$ over $K'$; it is fixed by $g^4$, and
$(L')^{\langle g^4\rangle}$ has degree $[L':K']/2=4$ over $K'$ by Galois theory. So they
coincide. The rest is the definition of $\res$ (Fact~\ref{fact:gc}(b)) and the
description of $T_2$ in Section~\ref{sec:facts}.
\end{proof}

\subsection{The witness curve and the descent field}

\begin{proposition}\label{prop:curve}
There exist a curve $C$ over $k$ with a faithful action of $G$ and a $G$-equivariant
$k$-algebra embedding $\iota\colon k(C)\hookrightarrow L'$.
\end{proposition}

\begin{proof}
By Fact~\ref{fact:geom} applied to $\tau_{\gen}$ and $F'=K'$ (with $m=1$; note
$T\times_KK'=\Spec L'$ is the spectrum of a field by Lemma~\ref{lem:Lprime}), there are
an irreducible quasi-projective $G$-variety $Z$ with faithful action, $\dim Z\le1$, and a
$G$-equivariant morphism $x\colon\Spec L'\to Z^{\rm free}$. If $\dim Z=0$ the converse
part of Fact~\ref{fact:geom} gives $\ed_k(\tau_{\gen}\otimes K')=0$, contradicting
Definition~\ref{def:witness}; so $Z$ is a curve. By Fact~\ref{fact:image}, $x$ is
dominant. Let $C$ be the smooth projective model of $Z$; the $G$-action on $Z$ extends to
$C$ (functoriality of normalisation and of the projective closure of a curve), faithfully,
and the dominant $x$ corresponds to a $k$-embedding $\iota\colon k(C)=k(Z)\hookrightarrow L'$,
equivariant because $x$ is.
\end{proof}

Fix $C$ and $\iota$ as in Proposition~\ref{prop:curve}. Since $G$ acts faithfully on
$C$, it acts faithfully on $k(C)$, and by Artin's theorem $k(C)/k(C)^G$ is Galois with
group $G$. Let $B:=C/G$ be the curve with $k(B)=k(C)^G$, and $t\colon C\to B$ the
quotient morphism (finite of degree $8$). As $\iota$ is equivariant,
$\iota(k(B))\subseteq(L')^G=K'$.

\begin{definition}\label{def:F0}
The \emph{descent field} of the witness is $F_0:=\iota(k(B))\subseteq K'$, a subfield
containing $k$ with $\tr_kF_0=1$. The \emph{descent class} is
$\beta\in W_3(F_0)/\wp W_3(F_0)=H^1(F_0,\Z/8)$, the class of the $G$-torsor
$\Spec k(C)\to\Spec k(B)$ transported along $\iota|_{k(B)}\colon k(B)\xrightarrow{\ \sim\ }F_0$.
\end{definition}

\begin{lemma}\label{lem:descentclass}
$s\equiv\beta$ in $W_3(K')/\wp W_3(K')$; that is, $\tau_{\gen}\otimes K'$ is the base
change to $K'$ of a $\Z/8$-torsor defined over $F_0$.
\end{lemma}

\begin{proof}
The $K'$-algebra $k(C)\otimes_{k(B)}K'$ (with $k(B)\to K'$ the map $\iota|_{k(B)}$) is
the base change of the $G$-torsor $\Spec k(C)\to\Spec k(B)$ along $\Spec K'\to\Spec k(B)$,
so it is a $G$-torsor over $K'$ of class $\res_{K'/F_0}(\beta)$. The map
$k(C)\otimes_{k(B)}K'\to L'$, $a\otimes c\mapsto\iota(a)c$, is a well-defined
$G$-equivariant $K'$-algebra homomorphism between $G$-torsors over $K'$, hence an
isomorphism (Fact~\ref{fact:gc}(b)). Since $L'/K'$ is $\tau_{\gen}\otimes K'$, of class
$\res_{K'/K}(s)$, the claim follows from Fact~\ref{fact:asw}.
\end{proof}

\subsection{The \texorpdfstring{$\Z/4$}{Z/4}-quotient and its twist}

Let $C_2:=C/\langle g^4\rangle$, the curve with $k(C_2)=k(C)^{\langle g^4\rangle}$.

\begin{lemma}\label{lem:C2}
$G_2=G/\langle g^4\rangle$ acts faithfully on $C_2$, with $C_2/G_2=B$, and
$\iota_2:=\iota|_{k(C_2)}\colon k(C_2)\hookrightarrow(L')^{\langle g^4\rangle}=L_2'$ is a
$G_2$-equivariant $k$-embedding. The corresponding morphism $x\colon\Spec L_2'\to C_2$
is dominant and satisfies $x\circ\Spec(g)=g_{C_2}\circ x$ for $g\in G_2$, where $\Spec(g)$
is the automorphism of $\Spec L_2'$ induced by the field automorphism $g$ of $L_2'$.
\end{lemma}

\begin{proof}
The elements of $G$ acting trivially on $k(C_2)=k(C)^{\langle g^4\rangle}$ form the
subgroup of $G=\Gal(k(C)/k(B))$ fixing this intermediate field pointwise, which is
$\langle g^4\rangle$ by the fundamental theorem of Galois theory. Hence $G_2$ acts
faithfully on $k(C_2)$, i.e.\ on $C_2$, and $k(C_2)^{G_2}=k(C)^G=k(B)$. The image of
$\iota_2$ lies in $(L')^{\langle g^4\rangle}=L_2'$ (Lemma~\ref{lem:Lprime}(b)) by
equivariance. A $k$-embedding of the function field of a curve into a field is the same
as a morphism from the spectrum of that field whose image is the generic point, and the
identity $\iota_2\circ g_{C_2}^*=g\circ\iota_2$ translates into
$x\circ\Spec(g)=g_{C_2}\circ x$.
\end{proof}

\begin{construction}[the twist $E_2$]\label{con:twist}
Let $Y:=C_2\times_k\Spec L_2$, an $L_2$-curve, and for $g\in G_2$ let
\[
\phi_g:=g_{C_2}\times\Spec(g)\colon Y\to Y ,
\]
where $\Spec(g)$ is induced by the field automorphism $g$ of $L_2$. Each $\phi_g$ is an
automorphism of the $k$-scheme $Y$ covering $\Spec(g)$ on $\Spec L_2$, and
$g\mapsto\phi_g$ is a group homomorphism because $G_2$ is abelian
($\Spec(g)\circ\Spec(h)=\Spec(hg)=\Spec(gh)$). So $(\phi_g)_g$ is a descent datum on
$Y$ relative to the Galois extension $L_2/K$ (whose group is $G_2$, Section~\ref{sec:facts}).
By Fact~\ref{fact:descent}(b) it is effective: there is a curve $E_2$ over $K$ with an
$L_2$-isomorphism
\[
\psi\colon E_2\otimes_KL_2\xrightarrow{\ \sim\ }Y=C_2\otimes_kL_2,\qquad
\psi\circ(\mathrm{id}_{E_2}\times\Spec g)=\phi_g\circ\psi\quad(g\in G_2).
\]
$E_2$ is the twist ${}^{(s_0,s_1)_K}C_2=(T_2\times_kC_2)/G_2$ of \cite[\S2.4]{P}
(the quotient of $Y=T_2\times_kC_2$ by the free action $\phi$), but we shall only use the
following consequences of the descent description, which agree with
Fact~\ref{fact:twistpts}.
\end{construction}

\begin{lemma}\label{lem:E2}
\begin{enumerate}[label=\textup{(\alph*)}]
\item For every field $F\supseteq K$ the map $P\mapsto x_P:=\mathrm{pr}_1\circ\psi\circ(P\times_K\mathrm{id}_{\Spec L_2})$
is a bijection from $E_2(F)$ onto the set of $k$-morphisms
$x\colon\Spec(F\otimes_KL_2)\to C_2$ with $x\circ\Spec(\mathrm{id}_F\otimes g)=g_{C_2}\circ x$ for all $g\in G_2$.
\item $K(E_2)=L_2(Y)^{G_2}$, where $L_2(Y)=\operatorname{Frac}(k(C_2)\otimes_kL_2)$ and
$g$ acts through $\phi_g^*$, i.e.\ by $g_{C_2}^*\otimes g$ on $k(C_2)\otimes_kL_2$.
Every element $h$ of $L_2(Y)^{G_2}$, viewed in $K(E_2)$, defines a $K$-morphism
$h\colon E_2\to\PP^1_K$ whose base change to $L_2$ is (via $\psi$) the morphism
$Y\to\PP^1_{L_2}$ defined by $h$.
\item There is a unique $K$-morphism $t_E\colon E_2\to B_K$ with
$(t\times\mathrm{id})\circ\psi=t_E\otimes\mathrm{id}_{L_2}$. For $P\in E_2(F)$,
$t_E(P)\in B_K(F)=B(F)$ is the unique $F$-point $b$ of $B$ with
$b\circ\mathrm{pr}=t\circ x_P$, where $\mathrm{pr}\colon\Spec(F\otimes_KL_2)\to\Spec F$.
\end{enumerate}
\end{lemma}

\begin{proof}
(a) By Fact~\ref{fact:descent}(a),(b), $\Hom_K(\Spec F,E_2)$ is the set of
$L_2$-morphisms $y\colon\Spec(F\otimes_KL_2)\to Y$ with
$\phi_g\circ y=y\circ\Spec(\mathrm{id}_F\otimes g)$ for all $g$, the bijection being
$P\mapsto\psi\circ(P\times\mathrm{id})$. An $L_2$-morphism $y$ to $Y=C_2\times_k\Spec L_2$
is the same as its first component $x=\mathrm{pr}_1\circ y$, a $k$-morphism to $C_2$
(the second component is the structure map). The compatibility condition has first
component $g_{C_2}\circ x=x\circ\Spec(\mathrm{id}\otimes g)$ and second component
$\Spec(g)\circ\mathrm{pr}_2\circ y=\mathrm{pr}_2\circ y\circ\Spec(\mathrm{id}\otimes g)$,
which holds automatically. (b) is Fact~\ref{fact:descent}(b); a rational function on a
curve over a field defines a morphism to $\PP^1$ compatibly with base change.
(c) $t\times\mathrm{id}\colon Y\to B\times_k\Spec L_2=B_K\otimes_KL_2$ satisfies
$(\mathrm{id}_B\times\Spec g)\circ(t\times\mathrm{id})=(t\times\mathrm{id})\circ\phi_g$
because $t\circ g_{C_2}=t$; by Fact~\ref{fact:descent}(b) it descends to a unique
$t_E$. For $P\in E_2(F)$, $(t_E(P))\times\mathrm{id}_{L_2}=(t\times\mathrm{id})\circ\psi\circ(P\times\mathrm{id})$
has first component $t\circ x_P$, and the first component of $b\times\mathrm{id}_{L_2}$
is $b\circ\mathrm{pr}$; finally $b\mapsto b\circ\mathrm{pr}$ is injective because
$\mathrm{pr}$ is faithfully flat, hence an epimorphism.
\end{proof}

\begin{definition}\label{def:Q}
$Q\in E_2(K')$ is the point with $x_Q=x$, the dominant equivariant morphism of
Lemma~\ref{lem:C2} (note $K'\otimes_KL_2=L_2'$, and $\Spec(\mathrm{id}_{K'}\otimes g)$
is the $\Spec(g)$ of Lemma~\ref{lem:C2}).
\end{definition}

\begin{lemma}[the descent point; $Q$ is non-constant]\label{lem:descentpt}
Let $b:=t_E(Q)\in B(K')$. Then the morphism $b\colon\Spec K'\to B$ has image the
generic point of $B$, and the induced map $b^*\colon k(B)\to K'$ equals
$\iota|_{k(B)}$; hence $b^*(k(B))=F_0$. In particular $b$ does not factor through any
$k$-point of $B$.
\end{lemma}

\begin{proof}
By Lemma~\ref{lem:E2}(c), $b\circ\mathrm{pr}=t\circ x$. The morphism $t\circ x$ has
image the generic point of $B$ ($x$ is dominant, $t$ surjective) and induces
$\iota\circ t^*=\iota|_{k(B)}$ on $k(B)$; hence $b$ has image the generic point, and
$\mathrm{pr}^*\circ b^*=\iota|_{k(B)}$ with $\mathrm{pr}^*$ the inclusion $K'\subseteq L_2'$.
\end{proof}

\begin{definition}[the involution and the divisor $D_Q$]\label{def:DQ}
$\Gal(K'/K)=\{1,\sigma\}$ acts on $E_2(K')$ by $P\mapsto\sigma P:=P\circ\Spec(\sigma)$.
Let $\xi_Q\in E_2$ be the image of $Q$, a closed point whose residue field $\kappa(\xi_Q)$
is embedded in $K'$ by $Q$, so $\deg\xi_Q:=[\kappa(\xi_Q):K]\in\{1,2\}$. Put
\[
D_Q:=\frac{2}{\deg\xi_Q}\,\xi_Q ,
\]
an effective $K$-rational divisor of degree $2$ on $E_2$, and
$\Delta:=\psi(D_Q\otimes_KL_2)$, an effective divisor of degree $2$ on $Y=C_2\otimes_kL_2$.
\end{definition}

\begin{lemma}\label{lem:DQ}
\begin{enumerate}[label=\textup{(\alph*)}]
\item $D_Q\otimes_KK'=Q+\sigma Q$ as divisors on $E_2\otimes_KK'$.
\item $\phi_g(\Delta)=\Delta$ for all $g\in G_2$.
\item $\Supp(D_Q\otimes_KL_2)=\xi_Q\times_K\Spec L_2$, and the morphism
$Q\times_K\mathrm{id}_{L_2}\colon\Spec L_2'\to E_2\otimes_KL_2$ factors through this closed
subscheme; its composite with $\mathrm{pr}_1\circ\psi$ is $x$.
\end{enumerate}
\end{lemma}

\begin{proof}
(a) If $\deg\xi_Q=2$ then $\kappa(\xi_Q)\cong K'$ and $\xi_Q\otimes_KK'=\Spec(K'\otimes_KK')$
consists of the two $K'$-points given by the two $K$-embeddings $\kappa(\xi_Q)\to K'$,
namely $Q$ and $\sigma Q$, which are distinct; so $D_Q\otimes K'=\xi_Q\otimes K'=Q+\sigma Q$.
If $\deg\xi_Q=1$ then $\xi_Q\otimes K'$ is the single $K'$-point $Q=\sigma Q$ and
$D_Q\otimes K'=2\xi_Q\otimes K'=Q+\sigma Q$. (b) $D_Q\otimes_KL_2$ is invariant under
$\mathrm{id}\times\Spec g$, and $\psi$ intertwines $\mathrm{id}\times\Spec g$ with
$\phi_g$. (c) $Q$ factors through $\xi_Q$, so $Q\times\mathrm{id}$ factors through
$\xi_Q\times_K\Spec L_2$, which is the support of $D_Q\otimes L_2$; the last assertion is
Definition~\ref{def:Q}.
\end{proof}

\section{Step 1: the class of \texorpdfstring{$D_Q$}{D\_Q} is constant}\label{sec:step1}

\begin{proposition}\label{prop:step1}
There is an effective divisor $D_0=P_1+P_2$ on $C_2$, with $P_1,P_2\in C_2(k)$, such that
\begin{enumerate}[label=\textup{(\alph*)}]
\item $\Delta\sim D_0\otimes_kL_2$ on $Y=C_2\otimes_kL_2$;
\item $g^*D_0\sim D_0$ on $C_2$ for every $g\in G_2$.
\end{enumerate}
\end{proposition}

\begin{proof}
\emph{(i) The class of $\Delta$ is constant.} By Fact~\ref{fact:pic}(b) (with $F=L_2$;
$C_2(k)\ne\emptyset$ as $k$ is algebraically closed) the class $[\Delta]\in\Pic(Y)$ is an
$L_2$-point $\lambda\colon\Spec L_2\to\Pic_{C_2/k}$, of degree $2$, i.e.\ a $k$-morphism
$\Spec L_2\to\Pic^2_{C_2/k}$. By Fact~\ref{fact:pic}(a), $\Pic^2_{C_2/k}$ is a torsor
under the Jacobian $J$ of $C_2$. The field $L_2=k(z_0,z_1,s_2)$ is purely transcendental
of transcendence degree $3$ over $k$: $K=k(s_0,s_1,s_2)\subseteq L_2\subseteq L$ with
$[L:K]=8$, so $\tr_kL_2=3$ and the three generators are algebraically independent.
Hence, by Fact~\ref{fact:rigid}(c), $\lambda$ factors through a $k$-point
$\lambda_0\in\Pic^2_{C_2/k}(k)=\Pic^2(C_2)$ (Fact~\ref{fact:pic}(b) with $F=k$), say
$\lambda_0=[D_1]$ with $D_1$ a divisor of degree $2$ on $C_2$. The identification of
Fact~\ref{fact:pic}(b) is natural in $F$, so the composite
$\Spec L_2\to\Spec k\xrightarrow{\lambda_0}\Pic^2_{C_2/k}$ is the class of
$D_1\otimes_kL_2$. Thus $\Delta\sim D_1\otimes_kL_2$ on $Y$.

\emph{(ii) The class contains an effective divisor.} Since $\Delta$ is effective,
$1\in L(\Delta)$, so $\ell(\Delta)\ge1$; linear equivalence preserves $\ell$, and
$\ell(D_1\otimes L_2)=\ell(D_1)$ by Fact~\ref{fact:bc}(a). So $\ell(D_1)\ge1$: choose
$0\ne h\in L(D_1)$ and put $D_0:=D_1+\divv h$, an effective divisor of degree $2$ on
$C_2$, linearly equivalent to $D_1$. Since $k$ is algebraically closed, $D_0=P_1+P_2$
with $P_i\in C_2(k)$. As $\divv(h)\otimes L_2=\divv(h)$, $D_0\otimes L_2\sim D_1\otimes L_2\sim\Delta$,
which is (a).

\emph{(iii) Invariance.} Let $g\in G_2$. Writing $D_0\otimes L_2=\mathrm{pr}_1^*D_0$ we get
\[
\phi_g^*(D_0\otimes L_2)=(\mathrm{pr}_1\circ\phi_g)^*D_0=(g_{C_2}\circ\mathrm{pr}_1)^*D_0
=\mathrm{pr}_1^*(g_{C_2}^*D_0)=(g^*D_0)\otimes L_2 .
\]
By Lemma~\ref{lem:DQ}(b), $\phi_g^*\Delta=\Delta$ (apply (b) to $g^{-1}$). Pullback along
the automorphism $\phi_g$ preserves linear equivalence, so (a) gives
$\Delta=\phi_g^*\Delta\sim\phi_g^*(D_0\otimes L_2)=(g^*D_0)\otimes L_2$, hence
$(g^*D_0)\otimes L_2\sim D_0\otimes L_2$ on $Y$. By Fact~\ref{fact:bc}(b) ($\Pic(C_2)\to\Pic(Y)$
is injective), $g^*D_0\sim D_0$ on $C_2$.
\end{proof}

\begin{remark}
The referee's formulation ``$\Pic_{E_2/K}$ is the twist of $\Pic_{C_2/k}$'' is not needed:
we only use $\Pic(Y)=\Pic_{C_2/k}(L_2)$ for the constant curve $C_2$, where the Picard
scheme is over $k$ and $C_2$ has $k$-points. The step ``$\ell(D_1)\ge1$, so $D_0$ may be
taken effective'' is the first of the two gaps flagged in the referee report; without it,
$\ell(D_0)\in\{1,2\}$ below would be unjustified (a general class of degree $2$ on a curve
of genus $\ge3$ has $\ell=0$).
\end{remark}

\section{Step 2: the dichotomy \texorpdfstring{$\ell(D_0)\in\{1,2\}$}{l(D0) in \{1,2\}}, and \texorpdfstring{$\ell(D_0)=1$}{l(D0)=1} is impossible}\label{sec:step2}

\begin{lemma}\label{lem:genus}
$g(C_2)\ge1$.
\end{lemma}

\begin{proof}
If $g(C_2)=0$ then $C_2\cong\PP^1_k$ and the faithful action of $G_2\cong\Z/4$ embeds
$\Z/4$ into $\Aut(\PP^1)=\mathrm{PGL}_2(k)$, contradicting Fact~\ref{fact:pgl}.
\end{proof}

\begin{corollary}\label{cor:l12}
$\ell(D_0)\in\{1,2\}$.
\end{corollary}

\begin{proof}
$D_0$ is effective, so $\ell(D_0)\ge1$; and $\ell(D_0)\le\deg D_0=2$ by
Fact~\ref{fact:rr}(b) and Lemma~\ref{lem:genus}.
\end{proof}

\begin{proposition}[$\ell(D_0)=1$ is impossible]\label{prop:step2a}
$\ell(D_0)=2$.
\end{proposition}

\begin{proof}
Suppose $\ell(D_0)=1$. By Fact~\ref{fact:bc}(a), $\ell(D_0\otimes L_2)=1$, and by
Proposition~\ref{prop:step1}(a), $\Delta=D_0\otimes L_2+\divv h$ with $0\ne h\in L(D_0\otimes L_2)=L_2\cdot1$;
so $h$ is a nonzero constant and
\[
\Delta=D_0\otimes_kL_2=(P_1\times\Spec L_2)+(P_2\times\Spec L_2)
\]
as divisors on $Y$, where $P_i\times\Spec L_2$ are the $L_2$-rational points of $Y$ over
$P_i$. In particular $\Supp\Delta\subseteq\{P_1\times\Spec L_2,P_2\times\Spec L_2\}$.
By Lemma~\ref{lem:DQ}(c), the morphism $\psi\circ(Q\times_K\mathrm{id}_{L_2})\colon\Spec L_2'\to Y$
factors through $\psi(\Supp(D_Q\otimes L_2))=\Supp\Delta$. Since $\Spec L_2'$ is a
one-point space, its image is one of the points $P_i\times\Spec L_2$, and composing with
$\mathrm{pr}_1\colon Y\to C_2$ we find that $x=\mathrm{pr}_1\circ\psi\circ(Q\times\mathrm{id})$
has image the $k$-point $P_i\in C_2$. This contradicts the dominance of $x$
(Lemma~\ref{lem:C2}). Hence $\ell(D_0)\ne1$, and Corollary~\ref{cor:l12} gives $\ell(D_0)=2$.
\end{proof}

\begin{remark}
This covers in particular every non-hyperelliptic $C_2$ of genus $\ge2$ (for which
$\ell(D)\le1$ for every divisor of degree $2$, by Fact~\ref{fact:rr}(c)), and every
hyperelliptic $C_2$ of genus $\ge2$ with $D_0\not\sim g^1_2$. The argument also covers the
case $Q\in E_2(K)$, where $D_Q=2\xi_Q$. Note that the contradiction is obtained from the
dominance of $x$, i.e.\ from the non-constancy of the descent point
(Lemma~\ref{lem:descentpt}): a constant $x$ would give a constant $b=t_E(Q)$.
\end{remark}

\begin{corollary}\label{cor:pencil}
There is a morphism $u\colon C_2\to\PP^1_k$ of degree $2$ such that $D_0\sim u^*(a)$ for
every $a\in\PP^1(k)$; $k(C_2)/k(u)$ is Galois of degree $2$, generated by an involution
$\iota_u$ of $C_2$ with $C_2/\iota_u\cong\PP^1$. If $g(C_2)\ge2$, then \emph{every}
degree-$2$ morphism $u'\colon C_2\to\PP^1$ satisfies $u'^*(a)\sim D_0$, and
$u'=\alpha\circ u$ for an automorphism $\alpha$ of $\PP^1$.
\end{corollary}

\begin{proof}
Fact~\ref{fact:rr}(c),(d) applied to $D_0$ (Proposition~\ref{prop:step2a},
Lemma~\ref{lem:genus}); for the last claim, $u'^*(a)$ has $\ell=2$ and degree $2$
(Fact~\ref{fact:rr}(c), converse part), so $u'^*(a)\sim D_0\sim u^*(a')$, and the proof of
Fact~\ref{fact:rr}(d) (case $e=2$) gives $u'=\alpha\circ u$.
\end{proof}

\section{Step 3: the hyperelliptic case, via a twisted pencil}\label{sec:step3}

\subsection{No descent field inside \texorpdfstring{$K$}{K}}

\begin{theorem}[{\cite[Theorem 11.6(b)]{A}}]\label{thm:caseA}
Let $K'=K(v_f)$ be a separable quadratic extension of $K$, $F_0\subseteq K'$ a subfield
containing $k$ with $\tr_kF_0=1$, and $\beta\in W_3(F_0)$ with $s\equiv\beta$ mod
$\wp W_3(K')$. Then $F_0\not\subseteq K$.
\end{theorem}

\begin{proof}
Suppose $F_0\subseteq K$. Then $s$ and $\beta$ both lie in $W_3(K)$ and have the same
image in $H^1(K',\Z/8)=W_3(K')/\wp$. Inflation--restriction for $K'/K$, with trivial
action on $\Z/8$, gives the exact sequence
\[
0\to H^1(\Gal(K'/K),\Z/8)\to H^1(K,\Z/8)\to H^1(K',\Z/8),
\]
with $H^1(\Gal(K'/K),\Z/8)=\Hom(\Z/2,\Z/8)\cong\Z/2$,
whose nonzero element is the image of the class $[f]\in H^1(K,\Z/2)$ of $K'/K$ under
the map induced by the inclusion $\Z/2\hookrightarrow\Z/8$, i.e.\ (Fact~\ref{fact:asw})
$V^2[f]=(0,0,f)$. Hence
\[
s\equiv\beta+\varepsilon\,(0,0,f)\pmod{\wp W_3(K)},\qquad\varepsilon\in\{0,1\}.
\]
Apply the truncation $W_3(K)\to W_2(K)$, a ring homomorphism commuting with $\wp$
(Fact~\ref{fact:asw}): $(0,0,f)$ maps to $0$, so
$(s_0,s_1)\equiv(\beta_0,\beta_1)\pmod{\wp W_2(K)}$ with $\beta_0,\beta_1\in F_0$.
Thus the $\Z/4$-torsor $(s_0,s_1)_K$ is the base change of a torsor over $F_0$, and
$\ed_k((s_0,s_1)_K)\le\tr_kF_0=1$. But $\ed_k((s_0,s_1)_K)=2$ by Fact~\ref{fact:ed}(b).
\end{proof}

\begin{corollary}\label{cor:F0notinK}
For our witness, $F_0=\iota(k(B))\not\subseteq K$.
\end{corollary}

\begin{proof}
Lemma~\ref{lem:descentclass} provides $\beta\in W_3(F_0)$ with $s\equiv\beta$ over $K'$.
\end{proof}

\subsection{The twisted pencil}

\begin{proposition}[pencil proposition]\label{prop:pencil}
Suppose there is $u\in k(C_2)$ with $[k(C_2):k(u)]=2$ such that
\begin{enumerate}[label=\textup{(H\arabic*)}]
\item $g^*u=u+1$ for the generator $g$ of $G_2$ (hence $(g^i)^*u=u+i$);
\item $D_0\sim u^*(a)$ for some (equivalently every) $a\in\PP^1(k)$.
\end{enumerate}
Then $F_0\subseteq K$, contradicting Corollary~\ref{cor:F0notinK}. In particular no such
$u$ exists.
\end{proposition}

\begin{proof}
\emph{(1) The quotient $B=\PP^1_t$.} Put $t:=u^2+u=q\circ u$ with $q\colon\PP^1_k\to\PP^1_k$,
$w\mapsto w^2+w$ (a morphism of degree $2$ with $q(\infty)=\infty$). By (H1), $t$ is
$G_2$-invariant, so $k(t)\subseteq k(C_2)^{G_2}=k(B)$; and
$[k(C_2):k(t)]=\deg(q\circ u)=\deg q\cdot\deg u=4=[k(C_2):k(B)]$. Hence $k(B)=k(t)$:
$B=\PP^1_t$, and the quotient morphism $t\colon C_2\to B$ is the composite $q\circ u$.

\emph{(2) The $K$-rational coordinate $\ut$.} Put
\[
\ut:=u\otimes1+1\otimes z_0\ \in\ k(C_2)\otimes_kL_2\subseteq L_2(Y).
\]
By (H1) and $g(z_0)=z_0+1$ (Section~\ref{sec:facts}),
$(g^*\otimes g)(\ut)=(u+1)\otimes1+1\otimes(z_0+1)=\ut$; so $\ut\in L_2(Y)^{G_2}=K(E_2)$
(Lemma~\ref{lem:E2}(b)). It defines a $K$-morphism $\ut\colon E_2\to\PP^1_K$ whose base
change to $L_2$ is, via $\psi$, the morphism $Y\to\PP^1_{L_2}$ given by $u+z_0$, that is
$\tau_{z_0}\circ u_{L_2}$ with $\tau_{z_0}\colon\PP^1_{L_2}\to\PP^1_{L_2}$, $w\mapsto w+z_0$.
For $c\in\PP^1(K)$ let $F_c:=\ut^*(c)$, an effective $K$-rational divisor on $E_2$
(Fact~\ref{fact:bc}(d)). By Fact~\ref{fact:bc}(d),
\[
\psi(F_c\otimes_KL_2)=(\tau_{z_0}\circ u_{L_2})^*(c)=u_{L_2}^*(\tau_{z_0}^{-1}c)=
\begin{cases}u_{L_2}^*(c+z_0)&c\in K,\\ u_{L_2}^*(\infty)&c=\infty,\end{cases}
\]
so $\deg F_c=\deg u=2$ and $\deg\ut=2$.

\emph{(3) Descent of the linear equivalence to $K$.} Any two points of $\PP^1(L_2)$ are
linearly equivalent on $\PP^1_{L_2}$, so all fibres $u_{L_2}^*(a')$, $a'\in\PP^1(L_2)$,
are linearly equivalent on $Y$; and $u_{L_2}^*(a)=u^*(a)\otimes_kL_2$ for $a\in\PP^1(k)$.
Hence, by (H2) and Proposition~\ref{prop:step1}(a),
\[
\psi(F_c\otimes L_2)\sim u^*(a)\otimes L_2\sim D_0\otimes L_2\sim\Delta=\psi(D_Q\otimes L_2)
\quad\text{on }Y,
\]
and transporting along the isomorphism $\psi^{-1}$: $D_Q\otimes_KL_2\sim F_c\otimes_KL_2$
on $E_2\otimes_KL_2$. By Fact~\ref{fact:bc}(b), $D_Q\sim F_c$ on $E_2$ for every
$c\in\PP^1(K)$. (This is the second gap flagged by the referee: the linear equivalence
obtained over $L_2$ does descend to $K$, because $\Pic(E_2)\to\Pic(E_2\otimes L_2)$ is
injective.)

\emph{(4) The $K$-rational members of $|F_c|$ are the $\ut$-fibres.} Take $c=0$ and write
$\Phi:=\ut^*(0)=\divv_0(\ut)$. By
Fact~\ref{fact:bc}(a), $\ell(\Phi)=\ell(\Phi\otimes L_2)=\ell(u^*(a)\otimes L_2)=\ell(D_0)=2$.
The functions $1$ and $\ut^{-1}$ lie in $L(\Phi)$ (indeed $\divv(\ut^{-1})+\Phi=\divv_\infty(\ut)\ge0$)
and are linearly independent over $K$ since $\ut$ is not constant; so
$L(\Phi)=K\cdot1\oplus K\cdot\ut^{-1}$. An effective divisor $D'$ on $E_2$ with $D'\sim\Phi$
is $D'=\Phi+\divv h$ with $0\ne h=\alpha+\beta\ut^{-1}$, $\alpha,\beta\in K$:
\begin{itemize}
\item if $\beta=0$, $D'=\Phi=\ut^*(0)$;
\item if $\alpha=0$, $D'=\divv_0(\ut)-\divv(\ut)=\divv_\infty(\ut)=\ut^*(\infty)$;
\item if $\alpha\beta\ne0$, $h=\alpha(\ut-c')/\ut$ with $c':=\beta/\alpha\in K$, and
$D'=\divv(\ut-c')-\divv(\ut)+\divv_0(\ut)=\divv(\ut-c')+\divv_\infty(\ut)=\divv_0(\ut-c')=\ut^*(c')$,
using $\divv_\infty(\ut-c')=\divv_\infty(\ut)$.
\end{itemize}
Since $D_Q$ is effective and $D_Q\sim\Phi$ by (3), $D_Q=\ut^*(c')$ for some $c'\in\PP^1(K)$.

\emph{(5) $\ut(Q)=c'$.} The point $Q\colon\Spec K'\to E_2$ factors through
$\xi_Q\subseteq\Supp D_Q=\Supp\ut^*(c')=\ut^{-1}(c')$. Hence $\ut\circ Q\colon\Spec K'\to\PP^1_K$
factors through the closed point $c'\in\PP^1_K$, whose residue field is $K$; a
$K'$-point of $\PP^1_K$ with image a $K$-rational closed point $c'$ is the point $c'$
itself. So $\ut(Q)=c'\in\PP^1(K)\subseteq\PP^1(K')$.

\emph{(6) $t_E=q_K\circ\ut$.} Let $q_K\colon\PP^1_K\to\PP^1_K$ be $w\mapsto w^2+w+s_0$,
$\infty\mapsto\infty$. By (1) and Lemma~\ref{lem:E2}(c), the base change of
$t_E\colon E_2\to B_K=\PP^1_K$ to $L_2$ is, via $\psi$, $t\times\mathrm{id}=q_{L_2}\circ u_{L_2}$;
the base change of $q_K\circ\ut$ is $q_{L_2}\circ\tau_{z_0}\circ u_{L_2}$, and
\[
q_{L_2}(\tau_{z_0}(w))=(w+z_0)^2+(w+z_0)+s_0=w^2+w+(z_0^2+z_0+s_0)=w^2+w
\]
because $z_0^2+z_0=s_0$ and the characteristic is $2$. So the two morphisms agree after
base change to $L_2$, hence are equal (Fact~\ref{fact:bc}(c)).

\emph{(7) Conclusion.} By (5) and (6), $b=t_E(Q)=q_K(\ut(Q))=q_K(c')\in\PP^1(K)$. By
Lemma~\ref{lem:descentpt}, $b$ has image the generic point of $B=\PP^1_k$; if $c'=\infty$
then $b=\infty$ is a $k$-point, which is excluded. So $c'\in K$ and $b$ is the $K$-point
$c'^2+c'+s_0$ of $\PP^1_K$, i.e.\ $b^*(t)=c'^2+c'+s_0\in K$. By Lemma~\ref{lem:descentpt},
$\iota(t)=b^*(t)$ and $F_0=\iota(k(B))=\iota(k(t))=k(\iota(t))\subseteq K$.
\end{proof}

\begin{remark}
The formula $\ut=u+z_0$ is the case $n=1$ of the coordinate $h=y-z$ of
\cite[Prop.~5.2(a)]{P}: a $K'$-point of the twist of $\PP^1_u$ (with $G_2$ acting through
$u\mapsto u+1$) is a $G_2$-equivariant map $\Spec L_2'\to\PP^1_u$, i.e.\ an element $u(Q)$
of $L_2'\cup\{\infty\}$ with $g(u(Q))=u(Q)+1$; then $u(Q)+z_0$ is $G_2$-invariant, i.e.\
lies in $K'$. Our argument works with the function $\ut$ on $E_2$ rather than with
points, so that the linear system $|F_c|$ can be used.
\end{remark}

\subsection{Verification of the hypotheses in the hyperelliptic case}

\begin{proposition}[Step 3]\label{prop:step3}
$g(C_2)\ge2$ is impossible.
\end{proposition}

\begin{proof}
Assume $g(C_2)\ge2$ and let $u\colon C_2\to\PP^1$ be the degree-$2$ map of
Corollary~\ref{cor:pencil}, with $D_0\sim u^*(a)$; write $L(D_0)=\langle h_0,h_1\rangle$,
$u=h_1/h_0$, so that $k(u)\subseteq k(C_2)$ is the subfield generated by the ratios of
elements of $L(D_0)$ (these ratios are exactly the M\"obius transforms $(\alpha+\beta u)/(\gamma+\delta u)$).

\emph{(i) $G_2$ preserves $k(u)$.} Let $h\in\Aut(C_2)$. Then $h^*D_0$ has degree $2$ and
$\ell(h^*D_0)=\ell(D_0)=2$ ($h^*$ is an automorphism of $k(C_2)$ mapping $L(D_0)$ onto
$L(h^*D_0)$), so $h^*D_0\sim D_0$ by Fact~\ref{fact:rr}(d): $h^*D_0=D_0+\divv h'$ and
$L(h^*D_0)=h'^{-1}L(D_0)$. A ratio of two elements of $L(D_0)$ is mapped by $h^*$ to a
ratio of two elements of $L(h^*D_0)$, which is a ratio of two elements of $L(D_0)$. Hence
$h^*(k(u))=k(u)$, and $h\mapsto\rho(h):=h^*|_{k(u)}$ is a homomorphism
$\Aut(C_2)\to\Aut_k(k(u))=\mathrm{PGL}_2(k)$ (the group $G_2$ is abelian, so the
contravariance of $h\mapsto h^*$ is harmless).

\emph{(ii) $\ker\rho|_{G_2}=\{1,g^2\}$ and $\rho(g)$ has order $2$.} The kernel of $\rho$
on $\Aut(C_2)$ consists of the automorphisms acting trivially on $k(u)$, i.e.\ of the
elements of $\Gal(k(C_2)/k(u))=\{1,\iota_u\}$ (Corollary~\ref{cor:pencil}). By
Fact~\ref{fact:pgl}, $\rho(g)$ has order $\le2$, so $g^2\in\ker\rho\cap G_2\subseteq\{1,\iota_u\}$;
as $g^2\ne1$, $g^2=\iota_u$ and $\ker\rho\cap G_2=\{1,g^2\}$; since $g\notin\ker\rho$,
$\rho(g)$ has order exactly $2$.

\emph{(iii) Normal form.} By Fact~\ref{fact:pgl}, $\rho(g)$ is conjugate in
$\mathrm{PGL}_2(k)$ to $u\mapsto u+1$; conjugating means replacing $u$ by a M\"obius
transform $u'=(\alpha+\beta u)/(\gamma+\delta u)$, which is again a degree-$2$ function
with $u'^*(a)\sim u^*(a')\sim D_0$ (the fibres of $u'$ are the fibres of $u$). After this
change of coordinate, $g^*u=u+1$.

So (H1) and (H2) of Proposition~\ref{prop:pencil} hold, and Proposition~\ref{prop:pencil}
gives a contradiction.
\end{proof}

\section{Step 4: genus one}\label{sec:step4}

In this section $g(C_2)=1$. We choose an origin and regard $C_2$ as an elliptic curve
$E$; $g\in\Aut(E)$ is an automorphism of order $4$ of the curve $E$.

\subsection{Faithful \texorpdfstring{$\Z/4$}{Z/4}-actions on elliptic curves in characteristic \texorpdfstring{$2$}{2}}

\begin{lemma}\label{lem:Z4}
Let $E$ be an elliptic curve over $k$ with origin $O$ and $\varphi\in\Aut(E)$ an
automorphism of order $4$; write $\varphi=t_P\circ\alpha$ with $P\in E(k)$,
$\alpha\in\Aut(E,O)$ (Fact~\ref{fact:ell}(a)).
\begin{enumerate}[label=\textup{(\alph*)}]
\item If $E$ is supersingular, then $\alpha$ has order $4$ and there is $R\in E(k)$ with
$t_R\circ\varphi\circ t_{-R}=\alpha$; in particular $\varphi$ fixes the point $-R$.
\item If $E$ is ordinary, then $\alpha=1$ and $\varphi=t_P$ with $P$ of order $4$.
\end{enumerate}
\end{lemma}

\begin{proof}
The projection $\Aut(E)\to\Aut(E,O)$ is a homomorphism, so $\alpha^4=1$. If $\alpha=1$,
$\varphi=t_P$ has order $\operatorname{ord}P$. If $\alpha^2=1\ne\alpha$, then
$(\alpha-1)(\alpha+1)=0$ in the domain $\End(E)$ gives $\alpha=-1$, and
$\varphi^2=t_P\alpha t_P\alpha=t_Pt_{\alpha P}\alpha^2=t_{P-P}=1$, so $\varphi$ does not
have order $4$. Hence either $\alpha=1$ and $\operatorname{ord}P=4$, or $\alpha$ has
order $4$.
(a) If $E$ is supersingular, $E(k)[2]=0$ (Fact~\ref{fact:ell}(b)), so $E(k)$ has no
element of order $4$ and $\alpha$ has order $4$. For $R\in E(k)$,
$t_R\varphi t_{-R}=t_Rt_P\alpha t_{-R}=t_{R+P}t_{-\alpha R}\alpha=t_{P+(1-\alpha)R}\alpha$
(Fact~\ref{fact:ell}(a)). Since $\alpha\ne1$, $1-\alpha$ is a nonzero isogeny, surjective
on $E(k)$ (Fact~\ref{fact:ell}(c)); choose $R$ with $(1-\alpha)R=-P$. Then
$t_R\varphi t_{-R}=\alpha$, and $\varphi(-R)=t_{-R}\alpha t_R(-R)=t_{-R}(\alpha(O))=-R$.
(b) If $E$ is ordinary, $j(E)\ne0$ and $\Aut(E,O)=\{\pm1\}$ (Fact~\ref{fact:ell}(b)), so
$\alpha$ has order $\le2$, hence $\alpha=1$ and $\operatorname{ord}P=4$.
\end{proof}

\begin{remark}
\cite[Lemma C.2]{A} (= \cite[Lemma 1.3]{ed2}) says more in the supersingular case: the
six elements of order $4$ of $\Aut(E,O)\cong\mathrm{SL}_2(\F_3)$ are conjugate, so
$(E,\varphi)$ is isomorphic to $(E_0,\text{Witt action})$ of \cite[Theorem 6.2]{P}, with
$E_0\colon y_1^2+y_1=y_0^3+y_0^2$ and $(y_0,y_1)\mapsto(y_0+1,y_1+y_0)$. We do not need
this identification: the abstract argument below only uses the fixed point of (a).
\end{remark}

\subsection{The supersingular case}

\begin{proposition}[Step 4(i)]\label{prop:step4i}
$g(C_2)=1$ with $C_2$ supersingular is impossible.
\end{proposition}

\begin{proof}
Let $E=C_2$. By Lemma~\ref{lem:Z4}(a) the automorphism $g$ of $E$ has a fixed point,
which we call $\infty$ and take as the origin; then $g=\alpha\in\Aut(E,\infty)$ is a
group automorphism of order $4$ (Fact~\ref{fact:ell}(a)).

\emph{(a) $E(k)^{\alpha}=0$.} In the domain $\End(E)$, $(\alpha^2-1)(\alpha^2+1)=\alpha^4-1=0$
and $\alpha^2\ne1$, so $\alpha^2=-1$. With $a:=\alpha+\hat\alpha\in\Z$ and $\deg\alpha=1$,
Fact~\ref{fact:ell}(c) gives $\alpha^2-a\alpha+1=0$, i.e.\ $a\alpha=0$, so $a=0$ because
$\End(E)$ is torsion-free and $\alpha\ne0$. Hence $\deg(1-\alpha)=1+\deg\alpha-a=2$, and
$\ker(1-\alpha)(k)=\{R\in E(k):\alpha R=R\}$ is a group of order $\le2$
(Fact~\ref{fact:ell}(c)), hence killed by $2$, hence contained in $E(k)[2]=0$
(Fact~\ref{fact:ell}(b)).

\emph{(b) $D_0\sim2\infty$.} By Proposition~\ref{prop:step1}(b), $g^*D_0\sim D_0$, and
$g^*\infty=\infty$ since $g$ fixes $\infty$; so the class $[D_0-2\infty]\in\Pic^0(E)$ is
$g^*$-invariant. Under the group isomorphism $E(k)\to\Pic^0(E)$, $R\mapsto[R-\infty]$
\cite[III.3.4]{Sil}, $g^*$ corresponds to $R\mapsto g^{-1}(R)=\alpha^{-1}(R)$ (the pullback
of the point divisor $R$ along the automorphism $g$ is the point $g^{-1}(R)$, and
$g^{-1}(\infty)=\infty$). Let $R\in E(k)$ with $[R-\infty]=[D_0-2\infty]$; then
$\alpha^{-1}R=R$, so $R\in E(k)^\alpha=0$ by (a), $R=\infty$, and $D_0\sim2\infty$.

\emph{(c) A coordinate with $g^*u=u+1$.} By Fact~\ref{fact:rr}(a), $\ell(2\infty)=2$ and
$\ell(\infty)=1$; let $L(2\infty)=\langle1,u\rangle$. Then $u$ has pole divisor exactly
$2\infty$, so $\deg u=[k(E):k(u)]=2$ and $u^*(\infty)=2\infty$. Since $g(\infty)=\infty$,
$g^*$ maps $L(2\infty)$ into itself ($\divv(g^*h)=g^*\divv h\ge-2g^*\infty=-2\infty$), so
$g^*u=au+b$ with $a\in k^\times$, $b\in k$. As $(g^*)^4=1$ on $L(2\infty)$, comparing the
coefficient of $u$ gives $a^4=1$, i.e.\ $(a+1)^4=0$, $a=1$. If $b=0$ then $u$ would be
$G_2$-invariant, i.e.\ $u\in k(B)$, whence $[k(E):k(u)]\ge[k(E):k(B)]=4$, a contradiction.
So $b\ne0$, and replacing $u$ by $u/b$ we get $g^*u=u+1$.

Now (H1) holds, and (H2) holds because $D_0\sim2\infty=u^*(\infty)$ by (b). Proposition~\ref{prop:pencil}
gives a contradiction.
\end{proof}

\begin{remark}
In the model $(E_0,\text{Witt action})$ one can take $u=y_0$ and $t=y_0^2+y_0$; this is
the picture of \cite[Theorem 6.2]{P}, and \cite[Corollary 11.7]{A} is the special case
$C=\bar U_3$ (where $C_2=E_0$ and $E_2=E'\otimes_{K_2}K$). The argument above is the
generalisation to an arbitrary $\Z/8$-curve $C$ with supersingular $C/\langle g^4\rangle$;
it needs neither \cite[Prop.~7.5]{P} nor the explicit equations.
\end{remark}

\subsection{The ordinary case: the corestriction argument}

\begin{proposition}[Step 4(ii)]\label{prop:step4ii}
$g(C_2)=1$ with $C_2$ ordinary is impossible.
\end{proposition}

\begin{proof}
Let $E=C_2$ with any origin $O$. By Lemma~\ref{lem:Z4}(b), $g=t_P$ with $P\in E(k)$ of
order $4$, so $G_2=\{t_{iP}\}$ acts by translations. By Fact~\ref{fact:ell}(d) there is a
finite \'etale isogeny $\phi\colon E\to E'$ of degree $4$ with kernel $\langle P\rangle$
and $E'\cong E/\langle t_P\rangle$; since $k(E/\langle t_P\rangle)=k(E)^{G_2}=k(B)$, we may
take $B=E'$ (an elliptic curve) and $t=\phi$. We have the exact sequence
\[
0\to\Z/4\xrightarrow{\ i\mapsto iP\ }E\xrightarrow{\ \phi\ }B\to0
\]
of Fact~\ref{fact:gc}(c), with connecting maps $\delta_F\colon B(F)\to H^1(F,\Z/4)$.

\emph{(a) $\delta_{K'}(b)=\res_{K'/K}(s_0,s_1)$, where $b=t_E(Q)\in B(K')$.} By
Lemma~\ref{lem:E2}(c), $\phi\circ x=b\circ\mathrm{pr}$ for $x=x_Q\colon\Spec L_2'\to E$.
So $x$ factors through a $K'$-morphism $x'\colon\Spec L_2'\to\phi^{-1}(b)=E\times_{B,b}\Spec K'$.
The scheme $\phi^{-1}(b)$ is the $\Z/4$-torsor of class $\delta_{K'}(b)$, with $1\in\Z/4$
acting by $t_P=g_E$; the scheme $\Spec L_2'$ is the $\Z/4$-torsor of class
$\res_{K'/K}(s_0,s_1)$, with $1$ acting by $\Spec(g)$ (Lemma~\ref{lem:Lprime}(b),
Fact~\ref{fact:asw}); and $x\circ\Spec(g)=g_E\circ x$ (Lemma~\ref{lem:C2}). Hence $x'$ is
a $\Z/4$-equivariant morphism of $\Z/4$-torsors over $K'$, so it is an isomorphism
(Fact~\ref{fact:gc}(b)), and the two classes coincide.

\emph{(b) $\Nm_{K'/K}(b)\in B(k)$ and $\delta_K$ vanishes on it.} $\Nm(b)=b+\sigma b$ is
$\sigma$-invariant, hence lies in $B(K)$ (Fact~\ref{fact:descent}(a)), and $B(K)=B(k)$ by
Fact~\ref{fact:rigid}(d) applied to the elliptic curve $B$ over $k$ and the purely
transcendental extension $K=k(s_0,s_1,s_2)$ of $k$. For $b_0\in B(k)$, the torsor
$\phi^{-1}(b_0)\times_k\Spec K$ is the base change of the $\Z/4$-torsor $\phi^{-1}(b_0)$
over $k$, which has a $k$-point ($\phi$ is surjective on $k$-points) and is therefore
trivial; so $\delta_K(b_0)=0$.

\emph{(c) Contradiction.} By Fact~\ref{fact:gc}(c),(a) and (a),(b):
\begin{align*}
0=\delta_K(\Nm_{K'/K}b)&=\cores_{K'/K}\,\delta_{K'}(b)=\cores_{K'/K}\res_{K'/K}(s_0,s_1)\\
&=2\cdot(s_0,s_1)\qquad\text{in }H^1(K,\Z/4)=W_2(K)/\wp .
\end{align*}
By Fact~\ref{fact:witt}(d), $2\cdot(s_0,s_1)=(0,s_0^2)$, and by Fact~\ref{fact:witt}(b),
$(0,s_0^2)-(0,s_0)=(0,s_0^2+s_0)=\wp(0,s_0)\in\wp W_2(K)$; so $2\cdot(s_0,s_1)\equiv(0,s_0)$,
which is nonzero in $W_2(K)/\wp$ by Lemma~\ref{lem:Vinj}. Contradiction.
\end{proof}

\begin{remark}\label{rem:iso}
(i) Step 4(ii) uses neither Step 1 nor the non-constancy of $b$: it shows that
$\res_{K'/K}(s_0,s_1)$ is not of the form $\delta_{K'}(b)$ for any $\Z/4$-isogeny
$E\to B$ of elliptic curves over $k$, any separable quadratic $K'/K$ and any $b\in B(K')$.
(ii) The identification $\Z/4\cong\ker\phi$, $1\mapsto P$, was chosen to match
$g=t_P$; another choice would change $\delta$ by an automorphism of $\Z/4$ acting on
$H^1$, which does not affect the conclusion $2\cdot\delta_{K'}(b)\ne0$.
\end{remark}

\section{Completeness of the case analysis, proof of the Theorem, corollaries}\label{sec:assembly}

\begin{proposition}[completeness]\label{prop:complete}
Let $K'$ be a witness and $C_2=C/\langle g^4\rangle$ as in Section~\ref{sec:setup}. Then
exactly one of the following holds, and each leads to a contradiction:
\begin{enumerate}[label=\textup{(\roman*)}]
\item $g(C_2)=0$: impossible by Lemma~\ref{lem:genus};
\item $g(C_2)=1$ and $C_2$ is supersingular: impossible by Proposition~\ref{prop:step4i};
\item $g(C_2)=1$ and $C_2$ is ordinary: impossible by Proposition~\ref{prop:step4ii};
\item $g(C_2)\ge2$: impossible by Proposition~\ref{prop:step3}.
\end{enumerate}
Hence no witness exists.
\end{proposition}

\begin{proof}
The genus of $C_2$ is a non-negative integer, and an elliptic curve in characteristic $2$
is either supersingular or ordinary; so the four cases are exhaustive and mutually
exclusive. In cases (ii) and (iv) the input is the effective divisor $D_0$ of
Proposition~\ref{prop:step1} with $\ell(D_0)=2$ (Proposition~\ref{prop:step2a}); in
case (iv) Corollary~\ref{cor:pencil} and Proposition~\ref{prop:step3} apply, in case (ii)
Proposition~\ref{prop:step4i}; case (iii) uses only Lemma~\ref{lem:Z4}(b) and the
corestriction argument. Each case ends in a contradiction with one of: Fact~\ref{fact:pgl} (case (i));
the dominance of $x$, i.e.\ the non-constancy of the descent point
(Proposition~\ref{prop:step2a}, which is used in cases (ii) and (iv)); Corollary~\ref{cor:F0notinK},
i.e.\ Theorem~\ref{thm:caseA} (cases (ii) and (iv), through Proposition~\ref{prop:pencil});
Lemma~\ref{lem:Vinj} (case (iii)).
\end{proof}

\begin{proof}[Proof of Theorem~\ref{thm:main}]
Let $F'/K$ be a field extension of degree $2$ and suppose $\ed_k(\tau_{\gen}\otimes F')\le1$.
By Lemma~\ref{lem:reduce}, $F'=K(v_f)$ is a witness with $f\notin\{0,s_0\}+\wp K$, and
Sections~\ref{sec:setup}--\ref{sec:step4} apply; Proposition~\ref{prop:complete} gives a
contradiction. Hence $\ed_k(\tau_{\gen}\otimes F')\ge2$ for every quadratic $F'/K$. For
$F'=K(v_2)$, Fact~\ref{fact:ed}(f) gives $\ed_k(\tau_{\gen}\otimes K(v_2))\le2$, so equality
holds there, and $\edd2_k(\tau_{\gen})=\min_{[F':K]\le2}\ed_k(\tau_{\gen}\otimes F')=2$
(the extension $F'=K$ contributes $\ed_k(\tau_{\gen})\ge2$ by Fact~\ref{fact:ed}(f)).
Finally $d_1(\tau_{\gen})=\min\{d:\edd d_k(\tau_{\gen})\le1\}$: since
$\edd1=\ed_k(\tau_{\gen})\ge2$ and $\edd2=2$, $d_1\ge3$; since $\edd4=1$
(Fact~\ref{fact:ed}(f)), $d_1\le4$.
\end{proof}

\begin{corollary}\label{cor:consequences}
\begin{enumerate}[label=\textup{(\alph*)}]
\item For every quasi-projective generically free primitive $G$-variety $Z$ over $k$ of
dimension $\le1$, the twist ${}^{\tau_{\gen}}Z$ has no free closed point of degree $1$ or
$2$: $D(Z)\cap\{1,2\}=\emptyset$ in the notation of \cite[Definition 4.3]{P}. In
particular this holds for every curve $Z$ with faithful $G$-action, and for
$Z=\bar U_3$ it recovers \cite[Corollary 11.7]{A}.
\item $\ed_k(\tau_{\gen}\otimes K')\ge2$ for every quadratic $K'/K$; the families of
\cite[Theorem 11.9]{A} and the extensions of \cite[Theorem 7.3]{P} are no longer special.
\item The profile at $(2,3)$ reads
$\edd1\in\{2,3\}$, $\edd2=2$, $\edd3\in\{1,2\}$, $\edd4=\dots=\edd7=1$, $\edd8=0$;
$d_2\in\{1,2\}$, $d_1\in\{3,4\}$, $d_0=8$. Whether $\edd3=1$ (a cubic witness) is
untouched by this note: Steps 1 and 4(ii) use the norm over a quadratic extension and the
$\sigma$-invariance of $D_Q$ in an essential way.
\end{enumerate}
\end{corollary}

\begin{proof}
(a) A free closed point of degree $d\le2$ on ${}^{\tau_{\gen}}Z$ with residue field $F'$
gives, by \cite[Theorem 4.2, (2)$\Rightarrow$(3)$\Rightarrow$(1)]{P}, $[F':K]=d$ and
$\ed_k(\tau_{\gen}\otimes F')\le1$, contradicting Theorem~\ref{thm:main}. (b) is
Theorem~\ref{thm:main}. (c) combines Theorem~\ref{thm:main} with \cite[Theorem 7.1]{P}
and \cite[Theorem 11.16]{A}.
\end{proof}

\subsection{The two ``windows'' of the external notes are closed}

The addendum \cite[Remark 11.12]{A} and the note \cite[Cor.~7.2]{S} list
two remaining shapes for a quadratic witness: (1) the descent class is the \'etale
$\Z/8$-isogeny torsor of an ordinary elliptic curve at a non-constant point (``isogeny
shape''), and (2) $f$ a polynomial in $s_2$ of degree $\ge8$ (a witness curve which is the
Artin--Schreier--Witt curve of a polynomial Witt vector). Both are excluded by
Theorem~\ref{thm:main}, since no witness exists at all. For (1) the corestriction
argument gives the exclusion directly and at level $3$, and shows that the ``necessary
condition'' of \cite[Prop.~7.1]{S} is vacuous:

\begin{proposition}[the isogeny shape is impossible]\label{prop:isogeny3}
Let $0\to\Z/8\to C\xrightarrow{\phi}C'\to0$ be an \'etale isogeny of elliptic curves over
$k$ with cyclic kernel of order $8$, let $K'/K$ be a separable quadratic extension, and
let $b\in C'(K')$. Then $\delta_{K'}(b)\ne\res_{K'/K}(s)$ in $H^1(K',\Z/8)$. In
particular $\tau_{\gen}\otimes K'$ is never of the isogeny shape, for any $K'$ and any
$b$ (constant or not), and the group $C'(K)/\phi(C(K))$ of \cite[Prop.~7.1]{S} is
$C'(k)/\phi(C(k))=0$.
\end{proposition}

\begin{proof}
Suppose $\delta_{K'}(b)=\res_{K'/K}(s)$. As in Proposition~\ref{prop:step4ii}(b),(c),
$\Nm_{K'/K}(b)\in C'(K)=C'(k)$ (Fact~\ref{fact:rigid}(d)), $\delta_K$ vanishes on
$C'(k)$, and Fact~\ref{fact:gc} gives
$0=\delta_K(\Nm b)=\cores\,\delta_{K'}(b)=\cores\res(s)=2s$. By Fact~\ref{fact:witt}(d),
$2s=(0,s_0^2,s_1^2)$, and
\[
(0,s_0^2,s_1^2)-(0,s_0,s_1)=V\bigl(F(s_0,s_1)-(s_0,s_1)\bigr)=V\wp(s_0,s_1)=\wp V(s_0,s_1)\in\wp W_3(K)
\]
(Fact~\ref{fact:witt}(c)); so $2s\equiv(0,s_0,s_1)=V(s_0,s_1)$, which is nonzero in
$W_3(K)/\wp$ by Lemma~\ref{lem:Vinj}. Contradiction. The last assertion:
$C'(K)=C'(k)$ and $\phi$ is surjective on $k$-points.
\end{proof}

\begin{remark}
(i) Thus the assertion in \cite[Prop.~7.1]{S} that ``$Q+\sigma Q$ has exact order $8$ in
$C'(K)/\phi(C(K))$'' cannot hold: its hypothesis is contradictory. (ii) For a witness
whose cover $C\to B$ is \'etale with $g(B)\ge2$ one could run the same norm argument
after pulling the cyclic cover back from a $\Z/8$-isogeny of $\operatorname{Jac}(B)$
along an Abel--Jacobi map; this is not needed, since such a $C$ has $g(C_2)\ge2$ and
falls under Step~3 (or Step~2).
\end{remark}

\section{What is cited without proof}\label{sec:inputs}

For a reader checking the argument line by line, here is the complete list of external
inputs; everything else in this note is proved above.

\subsection*{From \cite{P}}
\begin{itemize}
\item \S2.2: the Artin--Schreier--Witt isomorphism $H^1(F,\Z/2^n)=W_n(F)/\wp$ and its
compatibility with truncation, $V$, and base change (Fact~\ref{fact:asw}); the formulas
(2), (3), (4) (Fact~\ref{fact:witt}(a)--(c); (3) and its consequence (d) are re-derived,
(5) is reproved as (e)).
\item \S2.3: $L/K$ is Galois with group $G$, $L^{\langle g^{2^j}\rangle}=k(z_0,\dots,z_{j-1},s_j,\dots)$,
and $G/\langle g^4\rangle$ acts on $(z_0,z_1)$ by the Witt translation of $W_2$ with
$\wp(z_0,z_1)=(s_0,s_1)$ (used for $T_2$, Lemma~\ref{lem:Lprime}).
\item Lemma 2.1 ($\ed_k(\tau_{\gen})=\ed_k(\Z/8)\le3$); Fact 2.2(a) (twist points; we use
it only to identify our $E_2$ with the twist, Construction~\ref{con:twist}); Lemma 2.3 and
Theorem 4.2, (1)$\Rightarrow$(3) and (2)$\Rightarrow$(1) (Fact~\ref{fact:geom},
Corollary~\ref{cor:consequences}(a)); Lemma 2.4 (Fact~\ref{fact:image}); Lemma 2.5(ii)
(Fact~\ref{fact:pgl}, reproved); Proposition 3.2(c); Proposition 3.6; Proposition 3.9;
Lemma 4.1(i); Lemma 4.1(ii) (Fact~\ref{fact:rigid}(d), whose proof is quoted);
Proposition 6.1 ($\ed_k(\Z/4)=2$); Theorem 7.1 (upper bounds $\edd2\le2$ via $K(v_2)$,
$\edd4=1$, and $\ed_k(\tau_{\gen})\ge2$).
\end{itemize}

\subsection*{From \cite{A}}
Nothing is used without proof: Theorem 11.6(b) is reproved as Theorem~\ref{thm:caseA},
and Lemma C.2 is replaced by Lemma~\ref{lem:Z4}. Corollary 11.7, Proposition 11.8,
Theorems 11.9--11.10 and Corollary 11.11 are only referred to in remarks.

\subsection*{Standard references}
\begin{itemize}
\item Galois descent for quasi-projective schemes, descent of morphisms and of points,
$X(M)^\Gamma=X(F)$ \cite[6.1--6.2]{BLR}, \cite[V.20]{SeAG} (Fact~\ref{fact:descent}).
\item Existence and structure of the Picard scheme of a curve over $k$; $\Pic^d$ is a
torsor under the Jacobian; $\Pic(C_F)=\Pic_{C/k}(F)$ when $C(k)\ne\emptyset$
\cite[8.1, 8.2, 8.4]{BLR}, \cite{Kl} (Fact~\ref{fact:pic}).
\item Weil's extension theorem and the absence of rational curves on abelian varieties
\cite[Thm.~3.1, Cor.~3.8]{Mi} (Fact~\ref{fact:rigid}(a),(b)).
\item Flat base change for $H^0$, morphisms to $\PP^1$ and fibre divisors
\cite[II.6, III.9.3]{Ha}; Riemann--Roch \cite[IV.1.3]{Ha} (Fact~\ref{fact:bc}(a),(d),
Fact~\ref{fact:rr}(a)); arithmetic genus of a $(2,2)$-curve on $\PP^1\times\PP^1$ and
$g(\tilde\Gamma)\le p_a(\Gamma)$ (proof of Fact~\ref{fact:rr}(d)).
\item Elliptic curves \cite{Sil}: $\Aut(E)=E(k)\rtimes\Aut(E,O)$ (III.4.8); dual isogeny,
trace and degree (III.6.2--6.3); $\#\ker\le\deg$ (III.4.10); quotients by finite
subgroups (III.4.12); supersingularity in characteristic $2$ (V.3.1, Ex.~5.7);
$\Aut(E,O)$ for $j\ne0$ and $j=0$ (App.~A, Prop.~1.2); $E(k)\cong\Pic^0(E)$ (III.3.4)
(Facts~\ref{fact:ell}).
\item Group cohomology: restriction, corestriction as morphisms of $\delta$-functors,
$\cores\circ\res=[F':F]$, inflation--restriction \cite[Ch.~I, \S1.5]{NSW}, \cite[I.2, I.5]{Se};
torsors and $H^1$ \cite[I.5]{Se} (Fact~\ref{fact:gc}).
\item Artin's theorem on fixed fields; quotients of curves by finite groups; the smooth
projective model of a curve and functoriality of its automorphisms.
\end{itemize}

\begin{thebibliography}{BLR}
\bibitem[A]{A} Addendum to [P]: \emph{Section 11 (supplements at $(2,2)$ and $(2,3)$) and
Appendices A--F}, 27 September 2026 (files \texttt{body\_sec11.tex}, \texttt{body\_app.tex}).
\bibitem[BF]{BF} G.~Berhuy and G.~Favi, \emph{Essential dimension: a functorial point of
view (after A.~Merkurjev)}, Doc. Math. 8 (2003), 279--330.
\bibitem[BLR]{BLR} S.~Bosch, W.~L\"utkebohmert and M.~Raynaud, \emph{N\'eron models},
Ergeb. Math. Grenzgeb. (3) 21, Springer, 1990.
\bibitem[E]{ed2} \emph{Q1 at $(p,n)=(2,3)$: quadratic witnesses for $\ed^{[2]}(\tau_{\gen})=1$ --- a research note}, file \texttt{ed2\_analysis.md}, September 2026.
\bibitem[DR]{DR} A.~Duncan and Z.~Reichstein, \emph{Versality of algebraic group actions
and rational points on twisted varieties}, J. Algebraic Geom. 24 (2015), 499--530.
\bibitem[Ha]{Ha} R.~Hartshorne, \emph{Algebraic geometry}, GTM 52, Springer, 1977.
\bibitem[Kl]{Kl} S.~L.~Kleiman, \emph{The Picard scheme}, in: Fundamental algebraic
geometry, Math. Surveys Monogr. 123, AMS, 2005, 235--321.
\bibitem[Mi]{Mi} J.~S.~Milne, \emph{Abelian varieties}, in: Arithmetic geometry (Storrs
1984), Springer, 1986, 103--150.
\bibitem[NSW]{NSW} J.~Neukirch, A.~Schmidt and K.~Wingberg, \emph{Cohomology of number
fields}, 2nd ed., Grundlehren 323, Springer, 2008.
\bibitem[P]{P} \emph{The finite-degree profile of essential dimension of a
$\Z/p^n$-torsor in characteristic $p$}, manuscript, September 2026.
\bibitem[S]{S} \emph{The second slot of the finite-degree profile of the generic
$\Z/8$-torsor in characteristic 2: polynomial witnesses of degree $\le12$ do not exist},
note, September 2026.
\bibitem[Se]{Se} J.-P.~Serre, \emph{Galois cohomology}, Springer, 1997.
\bibitem[SeAG]{SeAG} J.-P.~Serre, \emph{Algebraic groups and class fields}, GTM 117,
Springer, 1988.
\bibitem[Sil]{Sil} J.~H.~Silverman, \emph{The arithmetic of elliptic curves}, 2nd ed.,
GTM 106, Springer, 2009.
\end{thebibliography}

\end{document}
